% 原 PDF 第71—96页；第Vbis讲引言及第1—2节。
\SourcePage{71}
\SGAchapter{Vbis}{上同调下降技术}{Techniques de descente cohomologique}
\begin{center}B. Saint-Donat\end{center}
% 原页的本讲目录由译本目录重建。
\textbf{引言}

1. 设 $X$ 是一个\Term{topological-space}{拓扑空间}，$\mathcal U=(U_i)_{i\in I}$ 是 $X$ 的一个\Term{covering}{覆盖}，假定它或者是\Term{open-covering}{开覆盖}，或者是\Term{locally-finite}{局部有限}的\Term{closed-covering}{闭覆盖}。若 $F$ 是 $X$ 上的\Term{abelian-sheaf}{阿贝尔层}，则由 $\mathcal U$ 定义的 Leray \Term{spectral-sequence}{谱序列} [[5] II (5.2.4) 和 (5.4.1)]
\[
\check H^p(\mathcal U,\mathcal H^q(F))\Longrightarrow H^*(X,F) \tag{1.1}
\]
可以如下描述：

\Term{covering}{覆盖} $\mathcal U$ 定义了一个关于 $F$ \Term{functoriality}{函子性}的“Čech 型”\Term{resolution}{解消} $\mathcal C^*(\mathcal U,F)$（同上 (5.2.1)）。另一方面，对每个 $F$ 都有一个关于 $F$ \Term{functoriality}{函子性}的“\Term{canonical-flasque-resolution}{典范松弛解消}” $C^*(F)$（同上 (4.3)）。采用这些记号，在 $\mathcal U$ 为\Term{open-covering}{开覆盖}的情形，\Term{spectral-sequence}{谱序列} (1.1) 由\Term{double-complex}{双复形}
\[
\Gamma(X,\mathcal C^*(\mathcal U,C^*(F))) \tag{1.2}
\]
得到。在 $\mathcal U$ 为\Term{locally-finite}{局部有限}\Term{closed-covering}{闭覆盖}的情形，考虑\Term{double-complex}{双复形}
\[
\Gamma(X,C^*(\mathcal C^*(\mathcal U,F))). \tag{1.3}
\]

\SourcePage{72}
2. 我们来寻找这些\Term{double-complex}{双复形}的统一描述。以 $X_\circ$ 表示各 $U_i$ 的\Term{disjoint-coproduct}{不交和}\Term{topological-space}{拓扑空间}，以 $X_n$（$n\geq0$）表示 $X_\circ$ 在 $X$ 上与自身作的第 $n+1$ 次迭代\Term{fiber-product}{纤维积}：
\[
X_n=\prod_{i_0\ldots i_n\in I} U_{i_0}\cap\cdots\cap U_{i_n}
=\prod_{\sigma\in\operatorname{Hom}([n],I)}\bigcap_0^n U_{\sigma(i)}. \tag{2.1}
\]
这些 $X_n$ 构成一个\Term{topological-space}{拓扑空间}的\Term{simplicial-system}{单纯系统}；若以 $j_n$ 表示 $X_n$ 到 $X$ 的\Term{projection}{投影}，则有
\[
\mathcal C^n(\mathcal U,F)=j_{n*}j_n^*(F). \tag{2.2}
\]
注意，\Term{canonical-flasque-resolution}{典范松弛解消}的构造与\Term{restriction}{限制}到\Term{open-subset}{开集}以及沿\Term{closed-immersion}{闭浸入}取\Term{direct-image}{直接像}可交换。因此：

(a) 若 $\mathcal U$ 是\Term{open-covering}{开覆盖}，
\[
\mathcal C^q(\mathcal U,C^n(F))=j_{q*}j_q^*C^p(F)=j_{q*}C^p(j_q^*F); \tag{2.3}
\]
(b) 若 $\mathcal U$ 是\Term{locally-finite}{局部有限}\Term{closed-covering}{闭覆盖}，
\[
C^p(\mathcal C^q(\mathcal U,F))=C^p(j_{q*}j_q^*F)=j_{q*}C^p(j_q^*F). \tag{2.4}
\]
由此可见，要得到 (1.2) 和 (1.3) 的统一描述，只须对每个 $q$ 在 $X_q$ 上取 $j_q^*(F)$ 的“\Term{canonical-flasque-resolution}{典范松弛解消}”，再对这一\Term{resolution}{解消}应用\Term{functor}{函子} $j_{q*}$。

3. 对于 $X$ 上任意一个\Term{topological-space}{拓扑空间}的\Term{simplicial-system}{单纯系统}，前面的描述仍然有意义：
\[
\begin{aligned}
\Delta^\circ&\longrightarrow\operatorname{Top}/X,\\
[n]&\longmapsto X_n,
\end{aligned}\tag{3.1}
\]
它不一定具有 (2.1) 的形式。不过，由 $F$ 余\Term{augmentation}{增广}的\Term{double-complex}{双复形}
\[
F\longrightarrow (j_{q*}C^p(j_q^*(F)))_{p,q} \tag{3.2}
\]
一般并不给出 $F$ 的\Term{resolution}{解消}。

本文致力于寻找使 (3.2) 给出 $F$ 的一个\Term{resolution}{解消}的充分条件。在这种情形，\Term{spectral-sequence}{谱序列} (1.1) 推广为一个称作“\Term{descent-spectral-sequence}{下降谱序列}”的\Term{spectral-sequence}{谱序列}
\[
\check H^p\bigl(H^q(X_p,j_p^*(F))\bigr)\longrightarrow H^{p+q}(X,F). \tag{3.3}
\]
在“常系数”的情形，Segal 曾用其他方法，为诸如 $K$ 理论这样的其他“\Term{cohomology}{上同调}理论”得到类似的\Term{spectral-sequence}{谱序列}（参见 [8]）。
\SourcePage{73}
Segal 在 CW \Term{complex}{复形}\Term{category}{范畴}中工作：他使用一个“\Term{geometric-realization}{几何实现}”\Term{functor}{函子}，它为一个\Term{topological-space}{拓扑空间}的\Term{semisimple}{半单}纯\Term{complex}{复形}配以一个新的\Term{topological-space}{拓扑空间}；这个新\Term{space}{空间}应当与一个\Term{simplicial-topos}{单纯拓扑斯}所关联的\Term{topos}{拓扑斯}相比较［参见 (1.2.12)］。

4. 在第5节，我们将通过如下构造说明所得判据：对 $\mathbf C$ 上的任一\Term{analytic-space}{解析空间} $X$，利用\Term{resolution-of-singularities}{奇点解消}，构造 $X$ 上一个\Term{nonsingular}{非奇异}\Term{analytic-space}{解析空间}的\Term{simplicial-system}{单纯系统}
\[
[n]\longmapsto X_n,
\]
使得 (3.2) 给出 $F$ 的一个\Term{resolution}{解消}。若取 $F$ 为\Term{constant-sheaf}{常层} $\mathbf C$，特别就得到一个\Term{spectral-sequence}{谱序列}
\[
H^q(X_p,\mathbf C)\longrightarrow H^{p+q}(X,\mathbf C), \tag{4.1}
\]
它用\Term{nonsingular}{非奇异}\Term{analytic-space}{解析空间}的复\Term{cohomology}{上同调}来表示 $X$ 的复\Term{cohomology}{上同调}。此外，若 $X$ 是\Term{projective}{射影}的，则可假定所有 $X_p$ 都是\Term{projective}{射影}的：这正是为 $X$ 得到某种“Hodge 理论”的关键因素（参见 [2]）。

5. 把\Term{sheaf}{层} $F$ 换成一个\Term{bounded-below}{下有界}的\Term{sheaf}{层}\Term{complex}{复形}，前面的构造可以原样推广。它们\Term{derivation}{导出}\Term{derived-category}{导出范畴}中的“\Term{localization}{局部化}”技术：

我们知道，对 (2.1) 给出的 $X_\circ$，\Term{arrow}{箭头}
\[
j_\circ^*:D^b(X)\longrightarrow D^b(X_\circ)
\]
一般并不\Term{faithful}{忠实}；我们将证明，有时利用涉及各 $X_n$、比 $j_\circ^*(K^\bullet)$（$K^\bullet\in D^+(X)$）更精确的数据，可以重建\Term{complex}{复形} $K^\bullet$。

所得陈述将在第 XVII 讲的附录中使用，以便把\Term{functor}{函子} $Rf_!$（$f$ 是\Term{scheme}{概形}之间\Term{separated}{分离}的\Term{finite-type}{有限型}\Term{morphism}{态射}）的定义推广到不假定 $f$ \Term{compactifiable}{可紧化}的情形。

在这一应用中，不可能只考虑\Term{topological-space}{拓扑空间}；这里须用\Term{scheme}{概形}的\Term{etale-site}{平展景}代替它们。另一方面，为了完成证明，不仅要考虑\Term{space}{空间}的\Term{simplicial-system}{单纯系统}，也必须考虑\Term{multisimplicial}{多重单纯}系统。这说明、解释了我们一开始所采用的极大一般性，或者说为它提供了理由。

\SGAsection{1}{预备知识}
\textbf{1.1. 记号。}
\SourcePage{74}
\textbf{1.1.1.} 下文始终以 $\mathcal U$ 表示一个满足 $\mathbf Z\in\mathcal U$ 的\Term{universe}{宇宙}；所考虑的所有\Term{topos}{拓扑斯}都是 $\mathcal U$-\Term{topos}{拓扑斯}。

设 $T,T'$ 是两个\Term{topos}{拓扑斯}：一个\Term{morphism}{态射} $\varphi:T\to T'$ 是一对\Term{functor}{函子} $\varphi_*:T\to T'$ 和 $\varphi^*:T'\to T$，连同一个\Term{adjunction}{伴随}
\[
\operatorname{Hom}_{T'}(\cdot,\varphi_*\cdot)
\xrightarrow{\sim}\operatorname{Hom}_T(\varphi^*\cdot,\cdot),
\]
并要求 $\varphi^*$ \Term{left-exact}{左正合}（即保持\Term{finite}{有限}\Term{projective-limit}{投射极限}）。

设 $(T,\mathcal O_T),(T',\mathcal O_{T'})$ 是两个\Term{ringed-topos}{赋环拓扑斯}：从 $(T,\mathcal O_T)$ 到 $(T',\mathcal O_{T'})$ 的\Term{morphism}{态射}是一个对 $(\varphi,\theta)$，其中 $\varphi:T\to T'$ 是\Term{topos}{拓扑斯}\Term{morphism}{态射}，$\theta:\mathcal O_{T'}\to\varphi_*(\mathcal O_T)$ 是\Term{ring-homomorphism}{环同态}。

\textbf{1.1.2.} 我们将经常使用 [SGA 1 VI] 中阐述的\Term{fibred-category}{纤维范畴}语言；读者也可以参阅 [4]。这里只固定一些记号：若 $E\to B$ 是一个\Term{fibred-functor}{纤维化函子}（相应地，\Term{cofibred-functor}{余纤维化函子}），对于 $B$ 中的\Term{morphism}{态射} $m:i\to j$，以 $m^*:E_j\to E_i$（相应地，$m_*:E_i\to E_j$）表示与之关联的“\Term{inverse-image}{逆像}”（相应地，“\Term{direct-image}{直接像}”）\Term{functor}{函子}；这些\Term{functor}{函子}中的每一个在唯一的\Term{functor}{函子}\Term{isomorphism}{同构}意义下确定。若 $\varphi:E\to E'$ 是一个 $B$-\Term{functor}{函子}，则对 $B$ 的每个\Term{object}{对象} $i$，以 $\varphi_i:E_i\to E'_i$ 表示 $\varphi$ 在 $E_i$ 上的\Term{restriction}{限制}。

\textbf{1.1.3.} 最后，以 $\Delta$ 表示如下\Term{category}{范畴}：$\Delta$ 的\Term{object}{对象}是\Term{ordered-set}{有序集} $[n]=\{0,1,\ldots,n\}$，其\Term{morphism}{态射}是所有的（广义）递增映射。以 $\Delta^+$（相应地，$\Delta^-$）\Term{representing-object}{表示对象}与 $\Delta$ 相同、\Term{arrow}{箭头}为 $\Delta$ 中的\Term{monomorphism}{单态射}（相应地，\Term{epimorphism}{满态射}）的\Term{category}{范畴}。我们将随需要自由使用有关 $\Delta$ 的经典记号［参见 [3] II 2］。

\textbf{1.2. $D$-\Term{topos}{拓扑斯}。}
\begin{SGAresult}{定义}{1.2.1}
设 $D$ 是一个\Term{category}{范畴}，$E\to D$ 是一个\Term{kan-fibration}{纤维化}且余\Term{kan-fibration}{纤维化}的\Term{functor}{函子}。若满足下列条件，就称 $E$ 在 $D$ 上是以\Term{topos}{拓扑斯}为\Term{fiber}{纤维}的\Term{bifibred-category}{双纤维范畴}，或称 $E$ 是一个 $D$-\Term{topos}{拓扑斯}：

(a) 对 $D$ 的每个\Term{object}{对象} $i$，\Term{fiber-category}{纤维范畴} $E_i$ 是一个\Term{topos}{拓扑斯}。

(b) 对 $D$ 中的每个\Term{morphism}{态射} $m:i\to j$，存在一个\Term{topos}{拓扑斯}\Term{morphism}{态射} $f:E_j\to E_i$，使得 $f_*=m^*$、$f^*=m_*$。
\end{SGAresult}
\textbf{注记。} 考虑到 (a)，条件 (b) 也可以表述为：\Term{functor}{函子} $m_*$ \Term{left-exact}{左正合}。

当 $D=\Delta$（相应地，$\Delta\times\Delta$）时，以\Term{simplicial-topos}{单纯拓扑斯}（相应地，双\Term{simplicial-topos}{单纯拓扑斯}）称呼一个 $\Delta$-\Term{topos}{拓扑斯}（相应地，$\Delta\times\Delta$-\Term{topos}{拓扑斯}）。

实际中，我们将通过下列考虑遇到 $D$-\Term{topos}{拓扑斯}：
\begin{SGAresult}{定义}{1.2.2}
设 $\mathscr E$ 是 $D$ 上一个\Term{kan-fibration}{纤维化}且余\Term{kan-fibration}{纤维化}的\Term{category}{范畴}。
\SourcePage{75}
若 $\mathscr E^\circ$ 是一个 $D^\circ$-\Term{topos}{拓扑斯}，就称 $\mathscr E$ 在 $D$ 上是以\Term{topos}{拓扑斯}的\Term{duality}{对偶}为\Term{fiber}{纤维}的\Term{bifibred-category}{双纤维范畴}。
\end{SGAresult}
\begin{SGAresult}{注记}{1.2.3}
明确地说，$\mathscr E$ 在 $D$ 上是以\Term{topos}{拓扑斯}的\Term{duality}{对偶}为\Term{fiber}{纤维}的\Term{bifibred-category}{双纤维范畴}，当且仅当满足下列条件：

(a) 对 $D$ 的每个\Term{object}{对象} $i$，\Term{fiber-category}{纤维范畴} $\mathscr E_i$ 的\Term{opposite-category}{对偶范畴}是一个\Term{topos}{拓扑斯}。

(b) 对 $D$ 中的每个\Term{morphism}{态射} $m:i\to j$，\Term{functor}{函子} $m^*$ \Term{right-exact}{右正合}。
\end{SGAresult}
\textbf{1.2.4.} 读者可在第4节找到以\Term{topos}{拓扑斯}的\Term{duality}{对偶}为\Term{fiber}{纤维}的\Term{bifibred-category}{双纤维范畴}的例子。

\textbf{1.2.5.} 设 $\mathscr E\to B$ 是 $B$ 上以\Term{topos}{拓扑斯}的\Term{duality}{对偶}为\Term{fiber}{纤维}的\Term{bifibred-category}{双纤维范畴}，$X:D^\circ\to B$ 是一个\Term{functor}{函子}。读者可自行验证
\[
(D^\circ\mathbin{\times}_B\mathscr E)^\circ
\]
是一个 $D$-\Term{topos}{拓扑斯}；我们将它记作 $\underline X$。

\textbf{1.2.6.} 我们常把一个\Term{functor}{函子} $X:D^\circ\to B$ 称为 $B$ 的一个 $D$-\Term{object}{对象}，并以 $X_i$ 表示此\Term{functor}{函子}在 $D$ 的\Term{object}{对象} $i$ 上的像；$B$ 的 $D$-\Term{object}{对象}构成一个\Term{category}{范畴}，记作 $D^\circ B$。若 $S$ 是 $B$ 的一个\Term{object}{对象}，则 $B/S$ 的一个 $D$-\Term{object}{对象}称为 $B$ 的一个向 $S$ \Term{augmentation}{增广}的 $D$-\Term{object}{对象}。给定 $B$ 的一个向 $S$ \Term{augmentation}{增广}的 $D$-\Term{object}{对象}，显然\Term{equivalence}{等价}于给定一个\Term{functor}{函子}\Term{morphism}{态射} $X\to C_S^D$，其中 $X$ 是 $B$ 的一个 $D$-\Term{object}{对象}，$C_S^D$ 是 $B$ 中取常值 $S$ 的 $D$-\Term{object}{对象}。

当 $D=\Delta$ 时，称为\Term{simplicial-object}{单纯对象}（相应地，\Term{augmentation}{增广}\Term{simplicial-object}{单纯对象}）。我们也将使用\Term{bisimplicial-object}{双单纯对象}（取 $D=\Delta\times\Delta$）。

\textbf{1.2.7.} 现在假定\Term{category}{范畴} $B$ 有\Term{finite}{有限}\Term{fiber-product}{纤维积}。设 $f:R\to S$ 是 $B$ 中的一个\Term{arrow}{箭头}：从 $\Delta^\circ\times(B/S)^\circ$ 到 $\operatorname{Ens}$ 的双\Term{functor}{函子}
\[
([n],X)\longmapsto\operatorname{Hom}_{\operatorname{Ens}}([n],\operatorname{Hom}_S(X,R))
\]
通过取 $X_n$ 为\Term{functor}{函子}
\[
Z\longmapsto\operatorname{Hom}_{\operatorname{Ens}}([n],\operatorname{Hom}_S(Z,R))
\]
的一个代表，定义一个\Term{functor}{函子} $[n]\mapsto X_{[n]}=X_n$；这里
\[
X_n\simeq\underbrace{R\times_S R\times_S\cdots\times_S R}_{n+1\text{ 次}}.
\]
我们将这样构造的向 $S$ \Term{augmentation}{增广}的\Term{semisimplicial-object}{半单纯对象}记作 $[R|_fS]$ 或 $[R|S]$。

最后，若 $X,X'$ 是 $B$（或 $B/S$）中的两个\Term{semisimplicial-object}{半单纯对象}，$u:X\to X'$ 是一个\Term{functor}{函子}\Term{morphism}{态射}，则出于技术原因，引入在\Term{category}{范畴} $\Delta^\circ B$ 中计算的\Term{object}{对象} $[X|_uX']$。它可以解释为 $B$ 的一个\Term{bisimplicial-object}{双单纯对象}，此时我们将它记作 $[[X|_uX']]$：事实上，有\Term{category}{范畴}的一个\Term{canonical}{典范}\Term{isomorphism}{同构}
\[
\Delta^\circ(\Delta^\circ B)\xrightarrow{\sim}(\Delta\times\Delta)^\circ B.
\]
现在回到一般的 $D$-\Term{topos}{拓扑斯}概念。

\SourcePage{76}
\textbf{1.2.8.} 设 $E$ 是一个 $D$-\Term{topos}{拓扑斯}；以 $(E)$ 表示\Term{category}{范畴} $\mathscr{H}\mathrm{om}_D(D,E)$。设 $f:D'\to D$ 是一个\Term{functor}{函子}：\Term{category}{范畴} $D'\times_D E$ 是一个 $D'$-\Term{topos}{拓扑斯}，与 $f$ 复合得到一个\Term{functor}{函子}
\[
f^*:(E)\longrightarrow(D'\times_D E).
\]
若 $D'$ 仅由 $D$ 的一个\Term{object}{对象} $i$（只有 $i$ 的恒等\Term{morphism}{态射}）组成，$f$ 取为记作 $e_i$ 的\Term{canonical}{典范}包含\Term{functor}{函子}，则可以取\Term{fiber-category}{纤维范畴} $E_i$ 作为 $D'\times_D E$，而 $e_i^*$ 就等同于“在 $i$ 处求值”\Term{functor}{函子}。

\begin{SGAresult}{命题}{1.2.9}
若 $D'$ 是一个 $\mathcal U$-\Term{small-category}{小范畴}，则\Term{functor}{函子} $f^*$ 有\Term{right-adjoint}{右伴随}和\Term{left-adjoint}{左伴随}（分别记作 $f_*,f_!$）。
\end{SGAresult}
这是 Kan 引理 [III 1.1] 的一个略微推广的结果；我们将经常使用它。

\begin{SGAresult}{引理}{1.2.10}
设 $I,J,A$ 是同一\Term{category}{范畴} $B$ 上的三个\Term{category}{范畴}；假定 $I$ 是 $\mathcal U$-小的，且 $A$ 在 $B$ 上\Term{kan-fibration}{纤维化}且余\Term{kan-fibration}{纤维化}。给定一个 $B$-\Term{functor}{函子} $f:I\to J$，以 $f^*$ 表示通过与 $f$ 复合定义的\Term{functor}{函子}
\[
\mathscr{H}\mathrm{om}_B(J,A)\longrightarrow\mathscr{H}\mathrm{om}_B(I,A).
\]
若 $A$ 在 $B$ 上的每个\Term{fiber}{纤维}中都存在 $\mathcal U$-\Term{inductive-limit}{归纳极限}（相应地，\Term{projective-limit}{投射极限}），则 $f^*$ 有\Term{left-adjoint}{左伴随}（相应地，\Term{right-adjoint}{右伴随}）。
\end{SGAresult}
\begin{proof}
我们只说明\Term{left-adjoint}{左伴随}\Term{functor}{函子}的存在性，引理中括号内的相应部分由\Term{duality}{对偶}性得到。将使用下面的事实，其验证留给读者：
\begin{SGAresult}{引理}{1.2.10.1}
设 $A$ 是\Term{category}{范畴} $B$ 上的\Term{bifibred-category}{双纤维范畴}，$b$ 是 $B$ 的\Term{object}{对象}，$(F_\lambda)_{\lambda\in\Lambda}$ 是从一个\Term{category}{范畴} $\Lambda$ 到\Term{fiber}{纤维} $A_b$ 的\Term{functor}{函子}。对 $A_{b'}$ 中的任意 $G$ 以及 $B$ 中的任意 $u:b'\to b$（相应地，$u:b\to b'$），只要左端有定义，就有\Term{bijection}{双射}
\[
\operatorname{Hom}_u(G,\varprojlim F_\lambda)\xrightarrow{\sim}\varprojlim\operatorname{Hom}_u(G,F_\lambda)
\]
（相应地，$\operatorname{Hom}_u(\varinjlim F_\lambda,G)\xrightarrow{\sim}\varprojlim\operatorname{Hom}_u(F_\lambda,G)$）。
\end{SGAresult}
以上准备就绪，令 $I\amalg_fJ$ 为由下列方式定义的 $B$ 上的\Term{category}{范畴}：
\[
\operatorname{ob}(I\amalg_fJ)=\operatorname{ob}(I)\amalg\operatorname{ob}(J),
\]
\[
\operatorname{Hom}(x,y)=
\begin{cases}
\operatorname{Hom}_I(x,y),&x,y\in\operatorname{ob}(I),\\
\operatorname{Hom}_J(x,y),&x,y\in\operatorname{ob}(J),\\
\operatorname{Hom}_J(f(x),y),&x\in\operatorname{ob}(I),\ y\in\operatorname{ob}(J),\\
\varnothing,&x\in\operatorname{ob}(J),\ y\in\operatorname{ob}(I).
\end{cases}
\]
\SourcePage{77}
\Term{category}{范畴} $\mathscr{H}\mathrm{om}_B(I\amalg_fJ,A)$ \Term{equivalence}{等价}于如下三元组的\Term{category}{范畴}：一个从 $I$ 到 $A$ 的 $B$-\Term{functor}{函子} $F$，一个从 $J$ 到 $A$ 的 $B$-\Term{functor}{函子} $G$，以及一个从 $F$ 到 $G\circ f=f^*(G)$ 的 $B$-\Term{functor}{函子}\Term{morphism}{态射} $\varphi$。

对 $J$ 的每个\Term{object}{对象} $j$，以 $I/j$ 表示 $I$ 中“由 $f$ 放在 $j$ 上方”的\Term{object}{对象}\Term{category}{范畴}：$I/j$ 的\Term{object}{对象}是对 $(i,\alpha)$，其中 $i$ 是 $I$ 的\Term{object}{对象}，$\alpha:i\to j$ 是 $I\amalg_fJ/j$ 中的\Term{arrow}{箭头}；$I/j$ 中的\Term{morphism}{态射}取为 $I\amalg_fJ/j$ 中的\Term{morphism}{态射}。若 $p_I,p_J$ 是 $I,J$ 到 $B$ 的\Term{projection}{投影}，且 $(i,\alpha)$ 是 $I/j$ 的\Term{object}{对象}，则以
\[
\alpha_*:A_{p_I(i)}\longrightarrow A_{p_J(j)}
\]
表示\Term{functor}{函子} $p_J(\alpha)_*$。

给定一个从 $I\amalg_fJ$ 到 $A$ 的 $B$-\Term{functor}{函子}，也就是给定 $F\in\operatorname{Hom}_B(I,A)$、$G\in\operatorname{Hom}_B(J,A)$ 和一个关于 $j$ \Term{functoriality}{函子性}的\Term{morphism}{态射}
\[
\psi_j:\varinjlim_{(i,\alpha)\in I/j}\alpha_*(F(i))\longrightarrow G(j).
\]
（左端关于 $j$ 的\Term{functoriality}{函子性}由 (1.2.10.1) 得到。）所以 $f^*$ 的\Term{left-adjoint}{左伴随}由下式给出：
\[
f_!(F)(j)=\varinjlim_{(i,\alpha)\in I/j}\alpha_*(F(i)).
\]
\end{proof}

\begin{SGAresult}{推论}{1.2.11}
设 $E$ 是一个 $D$-\Term{topos}{拓扑斯}，$i$ 是 $D$ 的\Term{object}{对象}；\Term{functor}{函子} $e_i^*$ 有一个\Term{left-adjoint}{左伴随}（相应地，\Term{right-adjoint}{右伴随}），由下式定义：
\[
e_{i!}(a)(j)=\coprod_{\alpha\in\operatorname{Hom}(i,j)}\alpha_*(a)
\quad\left(\text{相应地， }e_{i*}(a)(j)=\prod_{\alpha\in\operatorname{Hom}(j,i)}\alpha^*(a)\right),
\]
其中 $a$ 是 $E$ 中位于 $i$ 上方的\Term{object}{对象}。
\end{SGAresult}

\begin{SGAresult}{命题}{1.2.12}
设 $D$ 是一个 $\mathcal U$-\Term{small-category}{小范畴}，$E$ 是一个 $D$-\Term{topos}{拓扑斯}，则\Term{category}{范畴} $(E)$ 是一个 $\mathcal U$-\Term{topos}{拓扑斯}。
\end{SGAresult}
\begin{proof}
借助 (1.2.10.1)“逐\Term{fiber}{纤维}”验证，\Term{category}{范畴} $(E)$ 具有下列性质：

(a) \Term{finite}{有限}\Term{projective-limit}{投射极限}\Term{representable-functor}{可表示}。

(b) 由 $\mathcal U$ 的一个元素作指标的\Term{direct-sum}{直和}\Term{representable-functor}{可表示}，且是不交的、\Term{universal}{普适}的。

(c) \Term{equivalence-relation}{等价关系}是\Term{universal}{普适}\Term{effective}{有效的}。

还需证明 $(E)$ 有一个由 $\mathcal U$ 的一个元素作指标的\Term{generating-family}{生成元系}：若对 $D$ 的每个\Term{object}{对象} $i$，$(G_{i\lambda})_{\lambda\in\Lambda_i}$ 是 $E_i$ 的一个\Term{generating-family}{生成元系}（其中 $\Lambda_i$ 是一个 $\mathcal U$-\Term{small-set}{小集合}），则族 $(e_{i!}(G_{i\lambda}))_{i,\lambda}$ 是 $(E)$ 的一个\Term{generating-family}{生成元系}。
\end{proof}

\textbf{1.2.13.} 现在引入 $D$-\Term{topos}{拓扑斯}之间的\Term{morphism}{态射}概念。首先明确：若 $F,F'$ 是同一\Term{category}{范畴} $B$ 上的两个\Term{category}{范畴}，$T:F\to F'$ 是一个 $B$-\Term{functor}{函子}，则 $T$ 的一个 $B$-\Term{left-adjoint}{左伴随}是一个\Term{left-adjoint}{左伴随}于 $T$ 的\Term{functor}{函子} $S:F'\to F$，使\Term{canonical}{典范}\Term{morphism}{态射} $1\to T\circ S$
\SourcePage{78}
和 $S\circ T\to1$ 都是\Term{functor}{函子}的 $B$-\Term{morphism}{态射}。在这些条件下，显然可验证 $T$ 是\Term{cartesian}{笛卡尔}的，而 $S$ 是\Term{cocartesian}{余笛卡尔}的。

\begin{SGAresult}{定义}{1.2.14}
设 $E,E'$ 是两个 $D$-\Term{topos}{拓扑斯}：从 $E$ 到 $E'$ 的一个\Term{morphism}{态射}是一对 $D$-\Term{functor}{函子} $\Phi_*:E\to E'$ 和 $\Phi^*:E'\to E$，连同一个 $D$-\Term{adjunction}{伴随}
\[
\operatorname{Hom}(\Phi^*\cdot,\cdot)\xrightarrow[\sim]{\xi}(\cdot,\Phi_*\cdot),
\]
使得对 $D$ 的每个\Term{object}{对象} $i$，对 $(\Phi_{*i},\Phi_i^*)$ 连同 $\xi$ 诱导的\Term{adjunction}{伴随}，构成从 $E_i$ 到 $E'_i$ 的一个\Term{topos}{拓扑斯}\Term{morphism}{态射}。
\end{SGAresult}

\begin{SGAresult}{命题}{1.2.15}
设 $E,E'$ 是两个 $D$-\Term{topos}{拓扑斯}，$(\Phi_*,\Phi^*):E\to E'$ 是一个\Term{morphism}{态射}。假定 $D$ 是一个 $\mathcal U$-\Term{small-category}{小范畴}，则对
\[
(\Gamma(\Phi_*),\Gamma(\Phi^*)):(E)\longrightarrow(E')
\]
是一个\Term{topos}{拓扑斯}\Term{morphism}{态射}。
\end{SGAresult}
这显然由前一定义得到。

下面的引理可用来构造 $D$-\Term{topos}{拓扑斯}的\Term{morphism}{态射}，其证明留给读者。
\begin{SGAresult}{引理}{1.2.16}
设 $E,E'$ 是同一\Term{category}{范畴} $D$ 上的两个\Term{bifibred-category}{双纤维范畴}，$\Phi$ 是从 $E$ 到 $E'$ 的一个\Term{cartesian}{笛卡尔} $D$-\Term{functor}{函子}，并且对 $D$ 的每个\Term{object}{对象} $i$，$\Phi_i:E_i\to E'_i$ 有\Term{left-adjoint}{左伴随}。那么，对每个 $i$ 选择 $\Phi_i$ 的一个\Term{left-adjoint}{左伴随}，就\Term{canonical}{典范}地确定了一个 $D$-\Term{functor}{函子} $\Psi:E'\to E$，它是 $\Phi$ 的 $D$-\Term{left-adjoint}{左伴随}。
\end{SGAresult}
\begin{SGAresult}{附注}{1.2.17}
在 (1.2.16) 的条件下，假定 $E,E'$ 是两个 $D$-\Term{topos}{拓扑斯}，且对 $D$ 的每个\Term{object}{对象} $i$，\Term{functor}{函子} $\Phi_i$ 的任一\Term{left-adjoint}{左伴随}都是\Term{left-exact}{左正合}的。则若 $\psi:E'\to E$ 是 $\Phi$ 的一个 $D$-\Term{left-adjoint}{左伴随}，对 $(\Phi,\psi):E\to E'$ 就是一个 $D$-\Term{topos}{拓扑斯}\Term{morphism}{态射}。
\end{SGAresult}

现在设 $\mathscr E\to B$ 是 $B$ 上以\Term{topos}{拓扑斯}的\Term{duality}{对偶}为\Term{fiber}{纤维}的\Term{bifibred-category}{双纤维范畴}，$X,X'$ 是 $B$ 的两个 $D$-\Term{object}{对象}，$\alpha:X\to X'$ 是一个\Term{functor}{函子}\Term{morphism}{态射}。为 $\mathscr E$ 和 $\mathscr E^\circ$ 选择\Term{normalization}{正规化}\Term{cleavage}{劈裂}，就\Term{canonical}{典范}地确定了一个 $D$-\Term{topos}{拓扑斯}\Term{morphism}{态射}
\[
(\alpha_*,\alpha^*):\underline X\longrightarrow\underline{X'},
\]
使得对 $D$ 的每个\Term{object}{对象} $i$，
\[
(\alpha_{*i},\alpha_i^*):\underline X_i\longrightarrow\underline{X'}_i
\]
等于 $(\alpha_{i*}^\circ,\alpha_i^{*\circ})$，其中 $\alpha_i:X_i\to X'_i$ 是 $\alpha$ 给出的 $B$ 中的\Term{arrow}{箭头}。对于 $\mathscr E$ 和 $\mathscr E^\circ$ 的两个不同\Term{cleavage}{劈裂}选择，这样得到的\Term{morphism}{态射}之间存在唯一的 $D$-\Term{isomorphism}{同构}。（关于这些事实的验证，读者可以参阅 [4] (1.17)。）

\textbf{1.3. \Term{ringed}{赋环} $D$-\Term{topos}{拓扑斯}。}
\begin{SGAresult}{定义}{1.3.1}
\Term{ringed}{赋环} $D$-\Term{topos}{拓扑斯}是一个对 $(E,A)$，其中 $E$ 是 $D$-\Term{topos}{拓扑斯}，$A$ 是 $(E)$ 的一个环。
\end{SGAresult}
可验证，对 $D$ 的每个\Term{object}{对象} $i$，$A_i$ 是\Term{topos}{拓扑斯} $E_i$ 的一个环；对每个\Term{morphism}{态射} $m:i\to j$，\Term{canonical}{典范}\Term{arrow}{箭头} $A_i\to m^*(A_j)$ 是\Term{ring-homomorphism}{环同态}。

\SourcePage{79}
\begin{SGAresult}{定义}{1.3.2}
设 $(E,A),(E',A')$ 是两个\Term{ringed}{赋环} $D$-\Term{topos}{拓扑斯}：从 $(E,A)$ 到 $(E',A')$ 的一个\Term{morphism}{态射}是一个对 $(\Phi,\theta)$，其中 $\Phi:E\to E'$ 是 $D$-\Term{topos}{拓扑斯}\Term{morphism}{态射}，$\theta:A'\to\Gamma(\Phi_*(A))$ 是\Term{ring-homomorphism}{环同态}。
\end{SGAresult}
\begin{SGAresult}{注记}{1.3.3}
当 $D$ 是一个 $\mathcal U$-\Term{small-category}{小范畴}时（参见 1.2.15），\Term{ringed}{赋环} $D$-\Term{topos}{拓扑斯}的\Term{morphism}{态射} $\Phi:(E,A)\to(E',A')$ 诱导一个\Term{ringed}{赋环} $\mathcal U$-\Term{topos}{拓扑斯}的\Term{morphism}{态射}
\[
(\Gamma(\Phi),\theta):((E),A)\longrightarrow((E),A').
\]
\end{SGAresult}
\textbf{1.3.4.} 设 $\mathscr E\to B$ 是 $B$ 上以\Term{topos}{拓扑斯}的\Term{duality}{对偶}为\Term{fiber}{纤维}的\Term{bifibred-category}{双纤维范畴}，$\mathcal O$ 是 $\Gamma(\mathscr E^\circ)=\mathscr{H}\mathrm{om}_{B^\circ}(B^\circ,\mathscr E^\circ)$ 的一个环。若 $X:D^\circ\to B$ 是 $B$ 的一个 $D$-\Term{object}{对象}，则 $D$-\Term{topos}{拓扑斯} $\underline X$（参见 (1.2.5)）自然由 $(\mathcal O\circ X)^\circ:D\to\mathscr E$ \Term{ringed}{赋环}；以 $(\underline X,\mathcal O)$ 表示这样构造的\Term{ringed}{赋环} $D$-\Term{topos}{拓扑斯}。若 $\alpha:X\to X'$ 是一个\Term{functor}{函子}\Term{morphism}{态射}，则\Term{morphism}{态射} $(\alpha_*,\alpha^*):\underline X\to\underline{X'}$（参见 1.2.17）\Term{canonical}{典范}地诱导一个\Term{ringed}{赋环} $D$-\Term{topos}{拓扑斯}\Term{morphism}{态射} $(\underline X,\mathcal O)\to(\underline{X'},\mathcal O)$，仍记作 $\alpha$。

\textbf{1.3.5.} 一个\Term{ringed}{赋环} $D$-\Term{topos}{拓扑斯} $(E,A)$ \Term{canonical}{典范}地定义一个 $D$ 上以\Term{abelian-category}{阿贝尔范畴}为\Term{fiber}{纤维}的\Term{bifibred-category}{双纤维范畴} $\operatorname{Mod}(E,A)$，它在 $D$ 的\Term{object}{对象} $i$ 处的\Term{fiber}{纤维}，是\Term{ringed-topos}{赋环拓扑斯} $(E_i,A_i)$ 上的\Term{module-category}{模范畴} $\operatorname{Mod}(E_i,A_i)$。采用这些记号，$(E)$ 上的 $A$-\Term{module-category}{模范畴}记作 $\operatorname{Mod}((E),A)$，并与\Term{category}{范畴} $\mathscr{H}\mathrm{om}_D(D,\operatorname{Mod}(E,A))$ 等同。

\textbf{1.3.6.} 设 $\varphi=(\Phi,\theta):(E,A)\to(E',A')$ 是一个\Term{ringed}{赋环} $D$-\Term{topos}{拓扑斯}\Term{morphism}{态射}。它在相应的\Term{module-category}{模范畴}之间定义两个\Term{functor}{函子}
\[
\varphi_*:\operatorname{Mod}(E,A)\longrightarrow\operatorname{Mod}(E',A'),\qquad
\varphi^*:\operatorname{Mod}(E',A')\longrightarrow\operatorname{Mod}(E,A):
\]
设 $M$ 是 $\operatorname{Mod}(E,A)$ 中位于 $D$ 的\Term{object}{对象} $i$ 上方的\Term{object}{对象}：$\Phi_*(M)$ 是 $\Phi_*(A_i)$ 上的\Term{module}{模}，通过\Term{morphism}{态射} $\theta_i:A'_i\to\Phi_*(A_i)$，得到 $A'_i$ 上的一个\Term{module}{模}，记作 $\varphi_*(M)$。\Term{functor}{函子} $\varphi_*$ 称为沿\Term{morphism}{态射} $\varphi$ 的\Term{direct-image}{直接像}\Term{functor}{函子}。

设 $M'$ 是 $\operatorname{Mod}(E',A')$ 中位于 $D$ 的\Term{object}{对象} $i$ 上方的\Term{object}{对象}：$\Phi^*(M')$ 是 $\Phi^*(A'_i)$ 上的\Term{module}{模}，而
\[
\varphi^*(M')=\Phi^*(M')\otimes_{\Phi^*(A'_i)}A_i
\]
\Term{canonical}{典范}地具有 $A_i$-\Term{module}{模}结构。由 (1.2.16)，这样定义了一个\Term{left-adjoint}{左伴随}于 $\varphi_*$ 的\Term{functor}{函子} $\varphi^*$，称为沿\Term{morphism}{态射} $\varphi$ 的\Term{inverse-image}{逆像}\Term{functor}{函子}。

\textbf{1.3.6.1.} 若\Term{functor}{函子} $\Gamma(\varphi^*)$ \Term{exact}{正合}，就称 $\varphi$ \Term{flat}{平坦}。

\begin{SGAresult}{命题}{1.3.7}
设 $f:D'\to D$ 是一个\Term{functor}{函子}，$(E,A)$ 是一个\Term{ringed}{赋环} $D$-\Term{topos}{拓扑斯}。若 $D'$ 是一个 $\mathcal U$-\Term{small-category}{小范畴}，则\Term{canonical}{典范}\Term{functor}{函子}
\[
f^*:\operatorname{Mod}((E),A)\longrightarrow\operatorname{Mod}((E\times_D D'),A\circ f)
\]
有\Term{right-adjoint}{右伴随}和\Term{left-adjoint}{左伴随}。特别地，它是\Term{exact}{正合}的。
\end{SGAresult}
这立即由 (1.2.10) 和等同 $\mathscr{H}\mathrm{om}_D(D,\operatorname{Mod}(E,A))\simeq\operatorname{Mod}((E),A)$ 得到。按照一般记号，以 $f_!$（相应地，$f_*$）表示 $f^*$ 的\Term{left-adjoint}{左伴随}（相应地，\Term{right-adjoint}{右伴随}）。

\SourcePage{80}
\textbf{1.3.8.} 下面的讨论为下列类型的\Term{derived-functor}{导出函子}提供一种方便的计算方法：
\[
R^+\Gamma(\varphi_*):D^+((E),A)\longrightarrow D^+((E'),A'),
\]
其中 $\varphi:(E,A)\to(E',A')$ 是一个\Term{ringed}{赋环} $D$-\Term{topos}{拓扑斯}\Term{morphism}{态射}（参见 (1.3.6)）。［若 $(S,\mathcal O_S)$ 是\Term{ringed-topos}{赋环拓扑斯}，则以 $D^+(S,\mathcal O_S)$ 表示 $S$ 的 $\mathcal O_S$-\Term{module-category}{模范畴}的\Term{derived-category}{导出范畴}。］

\textbf{1.3.9.} 此外，这一形式计算也可用于其他背景，例如“$\ell$ 进\Term{cohomology}{上同调}的\Term{descent}{下降}”（参见 SGA 5）。

\begin{SGAresult}{命题}{1.3.10}
设 $D$ 是一个 $\mathcal U$-\Term{small-category}{小范畴}。若 $(E,A)$ 是\Term{ringed}{赋环} $D$-\Term{topos}{拓扑斯}，以 $I_{(E,A)}$ 表示 $\operatorname{Mod}((E),A)$ 中与形如
\[
\prod_{i\in\operatorname{ob}(D)}e_{i*}(Q_i)
\]
的\Term{object}{对象}\Term{isomorphism}{同构}的\Term{object}{对象}所成的\Term{set}{集合}，其中对 $D$ 的每个\Term{object}{对象} $i$，$Q_i$ 是\Term{totally-acyclic}{全无上同调}的（参见 V 4.1）\footnote{在第 V 讲的术语中称为“\Term{flasque}{松弛}”。}；$I_{(E,A)}$ 满足下列性质：

(i) 对 $\operatorname{Mod}((E),A)$ 的每个\Term{object}{对象} $F$ 以及 $D$ 的每个\Term{object}{对象} $i$，$e_i^*(F)$ 是\Term{totally-acyclic}{全无上同调}的。

(ii) $\operatorname{Mod}((E),A)$ 的每个\Term{object}{对象}都能\Term{injectivity}{单射}到 $I_{(E,A)}$ 的一个\Term{object}{对象}中。

(iii) $I_{(E,A)}$ 对\Term{finite}{有限}\Term{direct-sum}{直和}\Term{stable}{稳定}。

(iv) 对任一\Term{ringed}{赋环} $D$-\Term{topos}{拓扑斯}\Term{morphism}{态射} $\varphi:(E,A)\to(E',A')$，\Term{functor}{函子} $\Gamma(\varphi_*)$ 把 $C^+(\Gamma(E),A)$ 中由 $I_{(E,A)}$ 的\Term{object}{对象}组成的任一\Term{acyclic}{无上同调}\Term{complex}{复形}，变为一个由 $I_{(E',A')}$ 的\Term{object}{对象}组成的\Term{acyclic}{无上同调}\Term{complex}{复形}。
\end{SGAresult}
\begin{proof}
(i) 考虑到 $e_{i*}$ 的显式表达式（参见 (1.2.1.1)），只需证明下列引理：
\begin{SGAresult}{引理}{1.3.10.1}
设 $(S,\mathcal O_S)$ 是\Term{ringed-topos}{赋环拓扑斯}，$(F_t)_{t\in T}$ 是一个以 $\mathcal U$-\Term{small-set}{小集合} $T$ 为指标的\Term{totally-acyclic}{全无上同调} $\mathcal O_S$-\Term{module}{模}族。则 $\prod_{t\in T}F_t$ 是\Term{totally-acyclic}{全无上同调}的。
\end{SGAresult}
设 $X$ 是 $S$ 的\Term{object}{对象}：有一个\Term{spectral-sequence}{谱序列}
\[
H^p\left(X,\prod_{t\in T}^{(q)}F_t\right)
\longrightarrow\prod_{t\in T}H^{p+q}(X,F_t), \tag{1.3.10.2}
\]
其中 $\prod_{t\in T}^{(q)}$ 表示“以 $T$ 为指标的乘积”\Term{functor}{函子}的第 $q$ 个\Term{derived-functor}{导出函子}。由于 $\prod_{t\in T}^{(q)}F_t$ 是\Term{presheaf}{预层} $U\mapsto\prod_{t\in T}H^q(U,F_t)$ 的\Term{associated-sheaf}{伴随层}，(1.3.10.2) \Term{degeneration}{退化}，得到\Term{isomorphism}{同构}
\[
H^n\left(X,\prod_{t\in T}F_t\right)\xrightarrow{\sim}\prod_{t\in T}H^n(X,F_t), \tag{1.3.10.3}
\]
从而最终对每个 $n$ 都有 $H^n(X,\prod_{t\in T})=0$。

\SourcePage{81}
(ii) 设 $F$ 是 $\operatorname{Mod}((E),A)$ 的\Term{object}{对象}；对每个 $i$，选择一个\Term{monomorphism}{单态射} $F_i\to Q_i$，其中 $Q_i$ 在 $\operatorname{Mod}(E_i,A_i)$ 中\Term{totally-acyclic}{全无上同调}。则 $e_{i*}(F_i)\to e_{i*}(Q_i)$ 是\Term{monomorphism}{单态射}（因为 $e_{i*}$ 有\Term{left-adjoint}{左伴随}），而\Term{canonical}{典范}\Term{arrow}{箭头}
\[
F\longrightarrow\prod_{i\in\operatorname{Ob}(D)}e_{i*}(F_i)
\longrightarrow\prod_{i\in\operatorname{Ob}(D)}e_{i*}(Q_i) \tag{1.3.10.4}
\]
仍是\Term{monomorphism}{单态射}，这显然可以验证。

(iii) 证明留给读者。

(iv) \Term{functor}{函子} $\Gamma(\varphi_*)$ 把 $I_{(E,A)}$ 的\Term{object}{对象}变为 $I_{(E',A')}$ 的\Term{object}{对象}，这也留给读者验证。确定这一点后，考虑到 (i)，只需验证下列引理：
\begin{SGAresult}{引理}{1.3.10.5}
若 $0\to F\to G\to H\to0$ 是 $\operatorname{Mod}((E),A)$ 中的\Term{short-exact-sequence}{短正合列}，且 $F\in I_{(E,A)}$、$G\in I_{(E,A)}$，则列
\[
0\longrightarrow\Gamma(\varphi_*)(F)\longrightarrow\Gamma(\varphi_*)(G)
\longrightarrow\Gamma(\varphi_*)(H)\longrightarrow0
\]
在 $\operatorname{Mod}((E'),A')$ 中\Term{exact}{正合}。
\end{SGAresult}
只需注意，对 $D$ 的每个\Term{object}{对象} $i$，$e_i^*(H)$ 都\Term{acyclic}{无上同调}，而且 $\Gamma(\varphi_*)$ 的计算逐\Term{fiber}{纤维}进行。这就完成了 (1.3.10) 的证明。
\end{proof}

\begin{SGAresult}{推论}{1.3.11}
设 $\varphi:(E,A)\to(E',A')$ 是\Term{ringed}{赋环} $D$-\Term{topos}{拓扑斯}\Term{morphism}{态射}。可以使用由 $I_{(E,A)}$ 的\Term{object}{对象}组成的\Term{resolution}{解消}来计算\Term{functor}{函子} $R^+\Gamma(\varphi_*)$。
\end{SGAresult}
设 $L\subset K^+((E),A)$ 是由 $I_{(E,A)}$ 的\Term{object}{对象}组成的\Term{complex}{复形}所成的\Term{full-subcategory}{满子范畴}：由 ((1.3.10), (ii))，$L$ 是\Term{triangulated-subcategory}{三角子范畴}，可对它应用 [[7] 第一章] 的定理 (5.1)。

\begin{SGAresult}{推论}{1.3.12}
设 $\varphi:(E,A)\to(E',A')$ 是\Term{ringed}{赋环} $D$-\Term{topos}{拓扑斯}\Term{morphism}{态射}；对 $D$ 的每个\Term{object}{对象} $i$，以 $\varphi_i:(E_i,A_i)\to(E'_i,A'_i)$ 表示 $\varphi$ 在 $i$ 上方诱导的\Term{ringed}{赋环} $D$-\Term{topos}{拓扑斯}\Term{morphism}{态射}。则下图本质交换：
\[
\begin{tikzcd}
D^+(\Gamma(E),A) \arrow[r,"R^+(e_i^*)"] \arrow[d,"R^+\Gamma(\varphi_*)"'] & D^+(E_i,A_i) \arrow[d,"R^+\varphi_{i*}"]\\
D^+(\Gamma(E'),A') \arrow[r,"R^+(e_i^*)"'] & D^+(E'_i,A'_i).
\end{tikzcd}
\]
\end{SGAresult}
存在一个\Term{morphism}{态射}
\[
R^+(e_i^*)\circ R^+\Gamma(\varphi_*)\longrightarrow R^+\varphi_{i*}\circ R^+(e_i^*),
\]
由 (1.3.10) 可验证它是\Term{isomorphism}{同构}。

\SourcePage{82}
\begin{SGAresult}{注记}{1.3.13}
以 $a_{(E,A)}$ 表示 $\operatorname{Mod}((E),A)$ 中满足下列条件的\Term{object}{对象} $F$ 所成的\Term{set}{集合}：对 $D$ 的每个\Term{object}{对象} $i$，$e_i^*(F)$ 是\Term{totally-acyclic}{全无上同调}的。由 (1.3.10) 的证明，可以使用由 $a_{(E,A)}$ 的\Term{object}{对象}组成的\Term{resolution}{解消}计算 $R^+\Gamma(\varphi_*)$；特别地，引言中使用“\Term{canonical-flasque-resolution}{典范松弛解消}”\footnote{这一概念的推广参见第 XVII 讲。}时就是这样做的。由 [(1.3.10)(i)]，有 $I_{(E,A)}\subset a_{(E,A)}$，但在第2节我们将显式使用 $I_{(E,A)}$。
\end{SGAresult}

\SGAsection{2}{上同调下降的方法}
本节中，$D$ 是一个 $\mathcal U$-\Term{small-category}{小范畴}。

\textbf{2.1. 一般讨论。记号。}
\textbf{2.1.1.} 设 $(S,\mathcal O_S)$ 是\Term{ringed-topos}{赋环拓扑斯}。\Term{category}{范畴} $S\times D$ 配以\Term{canonical-projection}{典范投影} $S\times D\to D$，是一个 $D$-\Term{topos}{拓扑斯}；取常值 $\mathcal O_S$ 的\Term{section}{截面}进一步定义一个\Term{ringed}{赋环} $D$-\Term{topos}{拓扑斯} $(S\times D,\mathcal O_S)$，称为取常值 $(S,\mathcal O_S)$ 的\Term{ringed}{赋环} $D$-\Term{topos}{拓扑斯}。采用这些记号，\Term{category}{范畴} $\operatorname{Mod}((S\times D),\mathcal O_S)$ 等同于从 $D$ 到 $\operatorname{Mod}(S,\mathcal O_S)$ 的\Term{covariant-functor}{协变函子}\Term{category}{范畴}。

给 $S$ 上的每个\Term{module}{模} $F$ 配以取常值 $F$ 的\Term{functor}{函子}，便定义了一个\Term{exact}{正合}\Term{functor}{函子}
\[
\varepsilon^*:\operatorname{Mod}(S,\mathcal O_S)\longrightarrow\operatorname{Mod}((S\times D),\mathcal O_S).
\]
\Term{functor}{函子} $\varepsilon^*$ 有一个\Term{right-adjoint}{右伴随} $\varepsilon_*$，它给每个\Term{functor}{函子} $H:D\to\operatorname{Mod}(\mathcal O_S)$ 配以其\Term{projective-limit}{投射极限}；\Term{adjunction}{伴随}\Term{morphism}{态射} $F\to\varepsilon_*\varepsilon^*(F)$ 就是把 $F$ 映到取常值 $F$ 的\Term{functor}{函子}的\Term{projective-limit}{投射极限}中的\Term{morphism}{态射}。因此 $\varepsilon^*$ \Term{fully-faithful}{充分忠实}当且仅当 $D$ \Term{connected}{连通}。

\begin{SGAresult}{定义}{2.1.2}
设 $(E,A)$ 是\Term{ringed}{赋环} $D$-\Term{topos}{拓扑斯}；$(E,A)$ 的一个\Term{augmentation}{增广}是从 $(E,A)$ 到一个常值\Term{ringed}{赋环} $D$-\Term{topos}{拓扑斯}的\Term{morphism}{态射}（\Term{ringed}{赋环} $D$-\Term{topos}{拓扑斯}的\Term{morphism}{态射}）。配有一个\Term{augmentation}{增广}的\Term{ringed}{赋环} $D$-\Term{topos}{拓扑斯}称为\Term{augmentation}{增广}\Term{ringed}{赋环} $D$-\Term{topos}{拓扑斯}。
\end{SGAresult}
\textbf{2.1.3.} 设 $\mathscr E\to B$ 是以\Term{topos}{拓扑斯}的\Term{duality}{对偶}为\Term{fiber}{纤维}的\Term{bifibred-category}{双纤维范畴}，$\mathcal O$ 是 $(\mathscr E^\circ)=\mathscr{H}\mathrm{om}_{B^\circ}(B^\circ,\mathscr E^\circ)$ 的一个环。设 $S$ 是 $B$ 的\Term{object}{对象}；一个\Term{functor}{函子} $D^\circ\to B/S$，即一个\Term{functor}{函子}\Term{morphism}{态射} $X\xrightarrow{\theta}C_S^D$，其中 $X$ 是\Term{functor}{函子} $D^\circ\to B$（参见 (1.2.6)），诱导一个\Term{augmentation}{增广}
\[
\theta:(\underline X,\mathcal O)\longrightarrow(\mathscr E_S^\circ\times D,\mathcal O_S)
\qquad\text{（参见 (1.3.4)）}.
\]
我们将用产生它的\Term{functor}{函子}\Term{morphism}{态射}的同一字母表示这个\Term{augmentation}{增广}。

\textbf{2.2. \Term{cohomological-descent}{上同调下降}。}
\SourcePage{83}
\textbf{2.2.1.} 设 $\theta:(E,A)\to(S\times D,\mathcal O_S)$ 是一个\Term{augmentation}{增广}\Term{ringed}{赋环} $D$-\Term{topos}{拓扑斯}，采用 (1.3.6) 和 (2.1.1) 的记号，置
\[
\bar\theta^*=\Gamma(\theta^*)\circ\varepsilon^*:
\operatorname{Mod}(S,\mathcal O_S)\longrightarrow\operatorname{Mod}((E),A),
\]
以及 $\bar\theta^*=\varepsilon_*\circ\Gamma(\theta^*):\operatorname{Mod}((E),A)\to\operatorname{Mod}(S,\mathcal O_S)$。

$\bar\theta^*$ 的像落在 $\operatorname{Mod}((E),A)$ 的\Term{full-subcategory}{满子范畴}中，该\Term{full-subcategory}{满子范畴}由 $(D,\operatorname{Mod}(E,A))$ 的\Term{cocartesian}{余笛卡尔}\Term{section}{截面}组成；将这一\Term{category}{范畴}记作 $\operatorname{cocart}(D,\operatorname{Mod}(E,A))$。

\begin{SGAresult}{定义}{2.2.2}
若
\[
\Gamma(\theta^*)\circ\varepsilon^*:\operatorname{Mod}(S,\mathcal O_S)
\longrightarrow\operatorname{cocart}(D,\operatorname{Mod}(E,A))
\]
是\Term{category}{范畴}\Term{equivalence}{等价}，就称 $\theta$ 为\Term{effective-descent}{有效下降}\Term{augmentation}{增广}。
\end{SGAresult}
\textbf{2.2.3.} 采用 (2.2.1) 的记号，假定 $\theta$ \Term{flat}{平坦}，因而 $\bar\theta^*$ \Term{exact}{正合}，可直接\Term{descent}{下降}到\Term{derived-category}{导出范畴}；以 $L^+(\bar\theta^*)$ 表示所得\Term{functor}{函子}。采用这些记号：
\begin{SGAresult}{引理}{2.2.3.1}
存在两个\Term{functor}{函子}\Term{morphism}{态射}
\[
\begin{aligned}
\alpha &: \operatorname{id}_{D^+(S,\mathcal O_S)}\longrightarrow R^+(\bar\theta_*)\circ L^+(\bar\theta^*),\\
\beta &: L^+(\bar\theta^*)\circ R^+(\bar\theta_*)\longrightarrow\operatorname{id}_{D^+((E),A)},
\end{aligned}
\]
使这两个\Term{functor}{函子}构成\Term{adjunction}{伴随}。
\end{SGAresult}
（参见 [9] (3.3)。）

\begin{SGAresult}{定义}{2.2.4}
若满足下列条件，就称 $\theta$ 为 $1$-\Term{cohomological-descent}{上同调下降}\Term{augmentation}{增广}：

$1^\circ$）$\theta$ \Term{flat}{平坦}；

$2^\circ$）\Term{functor}{函子} $L^+(\bar\theta^*)$ \Term{fully-faithful}{充分忠实}。
\end{SGAresult}
\begin{SGAresult}{注记}{2.2.5}
前一定义的条件 $2^\circ$ 也可以表述为：(2.2.3.1) 中的\Term{morphism}{态射} $\alpha$ 是\Term{isomorphism}{同构}。
\end{SGAresult}
\begin{SGAresult}{定义}{2.2.6}
若 $\theta$ 同时是\Term{effective-descent}{有效下降}\Term{augmentation}{增广}和 $1$-\Term{cohomological-descent}{上同调下降}\Term{augmentation}{增广}，就称 $\theta$ 为 $2$-\Term{cohomological-descent}{上同调下降}\Term{augmentation}{增广}（或\Term{effective-cohomological-descent}{有效上同调下降}\Term{augmentation}{增广}）。
\end{SGAresult}
下述结果说明上述术语的理由：
\begin{SGAresult}{命题}{2.2.7}
设 $\theta$ 是一个\Term{effective-cohomological-descent}{有效上同调下降}\Term{augmentation}{增广}。则 $L^+(\bar\theta^*)$ 的\Term{essential-image}{本质像}是 $D^+((E),A)$ 的\Term{full-subcategory}{满子范畴}，由满足下列条件的\Term{complex}{复形} $F^\bullet$ 组成：对每个 $i$，$H^i(F^\bullet)$ 是 $\operatorname{Mod}(E,A)$ 的\Term{cocartesian}{余笛卡尔}\Term{section}{截面}。
\end{SGAresult}
\begin{proof}
为简便起见，将满足前述条件的\Term{complex}{复形} $F^\bullet$ 称为\Term{cohomological-descent-datum}{上同调下降数据}。显然，若 $K^\bullet\in D^+(S,\mathcal O_S)$，则 $L^+(\bar\theta^*)(K^\bullet)$
\SourcePage{84}
是\Term{cohomological-descent-datum}{上同调下降数据}；因此，只需证明对任一\Term{cohomological-descent-datum}{上同调下降数据} $F^\bullet$，\Term{canonical}{典范}\Term{morphism}{态射}
\[
\beta(F^\bullet):L^+(\bar\theta^*)\circ R^+(\bar\theta_*)(F^\bullet)\longrightarrow F^\bullet
\]
在 $D^+((E),A)$ 中是\Term{isomorphism}{同构}。

a) $F^\bullet$ 有界的情形：对满足 $H^i(F^\bullet)\ne0$ 的整数 $i$ 所在区间的长度 $\ell$ 作归纳：

当 $\ell\leq1$ 时，因为 $\theta$ 是\Term{effective-descent}{有效下降}\Term{augmentation}{增广}，断言成立。

假定 $\ell>1$，以 $n$ 表示满足 $H^n(F^\bullet)\ne0$ 的最大整数；有一个\Term{distinguished-triangle}{区别三角}\footnote{对每个\Term{complex}{复形} $K^\bullet$，$K^\bullet[-n]$ 表示\Term{complex}{复形} $T^{-n}(K^\bullet)$，其中 $T$ 是\Term{translation-functor}{平移函子}。}
\[
\begin{tikzcd}[column sep=large]
& H^n(F^\bullet)[-n]\arrow[dl] &\\
F'^\bullet\arrow[rr] &&F^\bullet\arrow[ul]
\end{tikzcd}
\]
使得归纳假设适用于 $F'^\bullet$。因为 $\beta(H^n(F^\bullet)[-n])$ 和 $\beta(F'^\bullet)$ 是\Term{isomorphism}{同构}，$\beta(F^\bullet)$ 也是\Term{isomorphism}{同构}。

b) 一般情形：以 $\sigma_{\leq}(F^\bullet)$ 表示\Term{complex}{复形}
\[
\cdots F^{n-1}\longrightarrow F^{n-1}\longrightarrow\operatorname{Ker}d^n
\longrightarrow0\longrightarrow0\cdots
\]
（参见 [7] (7.1)），因此对每个 $n$，有\Term{commutative-diagram}{交换图}
\[
\begin{tikzcd}[column sep=large]
L^+(\bar\theta^*)\circ R^+(\bar\theta_*)(\sigma_{\leq n}(F^\bullet)) \arrow[r,"\sim"] \arrow[d] & \sigma_{\leq n}(F^\bullet)\arrow[d]\\
L^+(\bar\theta^*)\circ R^+(\bar\theta_*)(F^\bullet)\arrow[r] & F^\bullet.
\end{tikzcd}
\]
只需证明\Term{morphism}{态射}
\[
R^+(\bar\theta_*)(\sigma_{\leq n}(F^\bullet))\longrightarrow R^+(\bar\theta_*)(F^\bullet)
\]
在 $i\leq n$ 时在 $H^i$ 上诱导\Term{isomorphism}{同构}。现在有一个三角：
\[
\begin{tikzcd}[column sep=large]
&F''^\bullet\arrow[dl]&\\
\sigma_{\leq n}(F^\bullet)\arrow[rr]&&F^\bullet\arrow[ul]
\end{tikzcd}
\]
\SourcePage{85}
其中 $i\leq n$ 时 $H^i(F''^\bullet)=0$。由此推出 $i\leq n$ 时 $H^i(R^+(\bar\theta_*)(F''^\bullet))=0$，再推出
\[
H^i(R^+(\bar\theta_*)(\sigma_{\leq n}(F^\bullet)))\longrightarrow H^i(R^+(\bar\theta_*)(F^\bullet))
\]
在 $i\leq n$ 时是\Term{isomorphism}{同构}，证明完成。
\end{proof}
现在介绍一种显式计算 $R^+\bar\theta_*$ 的方法，它对于证明某些\Term{cohomological-descent}{上同调下降}判据以及运用这一概念都很有用。

\textbf{2.3. $R^+\varepsilon_*$ 的一种计算方法。} 先给出两组一般性事实。

\textbf{2.3.1.} 设 $D$ 是一个 $\mathcal U$-\Term{small-category}{小范畴}，$\operatorname{Ab}$ 是属于 $\mathcal U$ 的\Term{abelian-group}{阿贝尔群}\Term{category}{范畴}。以 $r_D$ 表示由关系
\[
r_D(i)(j)=\mathbf Z^{\operatorname{Hom}(i,j)}
\]
定义的\Term{functor}{函子} $D\to(\mathscr{H}\mathrm{om}(D,\operatorname{Ab}))^\circ$，其中 $(i,j)$ 是 $D$ 的任一\Term{object}{对象}对。于是，对由 $D$ 的\Term{object}{对象} $i$ 和 $\mathscr{H}\mathrm{om}(D,\operatorname{Ab})$ 的\Term{object}{对象} $X$ 组成的每一对 $(i,X)$，构造一个\Term{canonical}{典范}且\Term{functoriality}{函子性}的\Term{isomorphism}{同构}
\[
\operatorname{Hom}(r_D(i),X)\xrightarrow{\sim}X(i).
\]
\textbf{2.3.1.1.} 设 $A$ 是\Term{additive-category}{加性范畴}（参见 Tohokû），$F$ 是\Term{functor}{函子} $D\to A$；通过关系
\[
\bar F(X)(a)=\operatorname{Hom}_{\mathscr{H}\mathrm{om}(D,\operatorname{Ab})}
\bigl(X,\operatorname{Hom}_A(a,F(\cdot))\bigr)
\]
定义\Term{functor}{函子} $\bar F:(\mathscr{H}\mathrm{om}(D,\operatorname{Ab}))^\circ\to A^\circ\operatorname{Ens}$，使得下图本质交换，其中 $h_A$ 表示\Term{canonical}{典范}\Term{functor}{函子}：
\[
\begin{tikzcd}
D\arrow[r,"r_D"]\arrow[d,"F"']&(\mathscr{H}\mathrm{om}(D,\operatorname{Ab}))^\circ\arrow[d,"\bar F"]\\
A\arrow[r,"h_A"']&\mathscr{H}\mathrm{om}(A^\circ,\operatorname{Ens}).
\end{tikzcd}
\]
以下留给读者验证：$\bar F$ 把 $\mathscr{H}\mathrm{om}(A,\operatorname{Ab})$ 中的\Term{finite}{有限}\Term{direct-sum}{直和}变为乘积；若以 $\mathbf Z$ 表示取常值 $\mathbf Z$ 的\Term{functor}{函子} $D\to\operatorname{Ab}$，则有 $\bar F(\mathbf Z)=\varprojlim F$。

最后，对应 $F\mapsto\bar F$ 定义一个\Term{covariant-functor}{协变函子}
\[
\mathscr{H}\mathrm{om}(D,A)\longrightarrow
\mathscr{H}\mathrm{om}\bigl((\mathscr{H}\mathrm{om}(A,\operatorname{Ab}))^\circ,
\mathscr{H}\mathrm{om}(A^\circ,\operatorname{Ens})\bigr).
\]
\textbf{2.3.1.2.} 以 $\operatorname{Add}(D)$ 表示 $(\mathscr{H}\mathrm{om}(D,\operatorname{Ab}))^\circ$ 的\Term{full-subcategory}{满子范畴}，由形如 $r_D(i)$（$i$ 为 $D$ 的\Term{object}{对象}）的\Term{object}{对象}及其\Term{finite}{有限}\Term{direct-sum}{直和}组成。\Term{category}{范畴} $\operatorname{Add}(D)$ 是\Term{additive}{加性}的，\Term{functor}{函子} $r_D:D\to\operatorname{Add}(D)$ 满足以下\Term{universal-property}{泛性质}：

\SourcePage{86}
\textbf{2.3.1.3.} 对每个\Term{additive-category}{加性范畴} $A$，\Term{functor}{函子} $G\mapsto G\circ r_D$ 从 $\mathscr{H}\mathrm{om}(\operatorname{Add}(D),A)$ 中由\Term{additive-functor}{加性函子}组成的\Term{full-subcategory}{满子范畴}到\Term{category}{范畴} $\mathscr{H}\mathrm{om}(D,A)$ 诱导一个\Term{equivalence}{等价}。

在\Term{equivalence}{等价}意义下由 (2.3.1.3) 确定的\Term{category}{范畴} $\operatorname{Add}(D)$，称为由 $D$ 生成的\Term{additive-category}{加性范畴}。

\textbf{2.3.2.} 设 $A$ 是\Term{abelian-category}{阿贝尔范畴}，$k^{\bullet\bullet}$ 是一个\Term{double-complex}{双复形}，我们把它视作 $C^+(C^+(A))$ 的\Term{object}{对象}，第一个指标对应第一个\Term{signature}{符号} $C^+$；以 $(K^{\bullet\bullet})_s$ 表示关联的\Term{simple-complex}{单复形}（参见 Tohokû (2.4)，并沿用其中的记号）。这样便定义了一个\Term{functor}{函子}
\[
(\ )_s:C^+(C^+(A))\longrightarrow C^+(A). \tag{2.3.2.1}
\]
下列引理的验证留给读者：
\begin{SGAresult}{引理}{2.3.2.2}
\Term{functor}{函子} $(\ )_s$ 定义一个保持\Term{quasi-isomorphism}{拟同构}的\Term{triangulated-functor}{三角函子}
\[
K^+(C^+(A))\longrightarrow K^+(A);
\]
因而定义一个\Term{triangulated-functor}{三角函子}
\[
D^+(C^+(A))\longrightarrow D^+(A),
\]
仍记作 $(\ )_s$。
\end{SGAresult}
\textbf{注记。} 对\Term{functor}{函子} $(\ )_s:C(C^b(A))\to C(A)$ 也有相同结果，因为由 (EGA III (11.3.3))，所考虑的\Term{spectral-sequence}{谱序列}是\Term{biregular}{双正则}的。

\textbf{2.3.3.} 以上准备就绪，设 $(S,\mathcal O_S)$ 是\Term{ringed-topos}{赋环拓扑斯}，$D$ 是 $\mathcal U$-\Term{small-category}{小范畴}。设 $Z^\bullet\in C^+(\operatorname{Add}(D))$ 是一个\Term{cochain-complex}{上链复形}，且 $n<0$ 时 $Z^n=0$。

由 $\operatorname{Add}(D)$ 的\Term{universal-property}{泛性质}，每个 $\operatorname{Mod}((S\times D,\mathcal O_S),\mathcal O_S)$ 的\Term{object}{对象} $F$（由 (2.1.1)，它等同于一个\Term{functor}{函子} $D\to\operatorname{Mod}(S,\mathcal O_S)$）都定义了 $C^+(S,\mathcal O_S)$ 的一个\Term{object}{对象} $\varepsilon_{Z^\bullet*}(F)$；由 (2.3.1.3)，它随 $F$ \Term{functoriality}{函子性}地变化。

这样便定义了一个\Term{triangulated-functor}{三角函子}
\[
K^+(\varepsilon_{Z^\bullet*}):K^+((S\times D),\mathcal O_S)
\longrightarrow K^+(C^+(S,\mathcal O_S)). \tag{2.3.3.1}
\]
此外，由 (2.3.2.2)，复合\Term{functor}{函子}
\[
(\ )_s\circ K^+(\varepsilon_{Z^\bullet*}):K^+((S\times D),\mathcal O_S)
\longrightarrow K^+(S,\mathcal O_S)
\]
把\Term{acyclic}{无上同调}\Term{object}{对象}变为\Term{acyclic}{无上同调}\Term{object}{对象}。因此它定义一个\Term{triangulated-functor}{三角函子}
\[
R^+\varepsilon_{Z^\bullet*}:D^+((S\times D),\mathcal O_S)
\longrightarrow D^+(S,\mathcal O_S), \tag{2.3.3.2}
\]
满足关系 $Q\circ(\ )_s\circ K^+(\varepsilon_{Z^\bullet*})=R^+\varepsilon_{Z^\bullet*}\circ Q$。\footnote{对任一\Term{abelian-category}{阿贝尔范畴} $A$，字母 $Q$ 表示\Term{canonical}{典范}\Term{functor}{函子} $K^+(A)\to D^+(A)$。}

\SourcePage{87}
\textbf{2.3.4.} 保持前述数据不变。将 $Z^\bullet$ 看作 $\mathscr{H}\mathrm{om}(D,\operatorname{Ab})$ 中的一个\Term{chain-complex}{链复形}。设 $t:Z^\circ\to\mathbf Z$ 是从 $Z^\circ$ 到取常值 $\mathbf Z$ 的\Term{functor}{函子}的一个\Term{morphism}{态射}，使复合\Term{morphism}{态射}
\[
Z^1\longrightarrow Z^0\xrightarrow{t}\mathbf Z \tag{2.3.4.1}
\]
为零。

由\Term{functoriality}{函子性}（参见 (2.3.1.1)），得到一个\Term{morphism}{态射} $\varepsilon_*(F)\to\varepsilon_{Z^\bullet*}(F)^\circ$，使复合\Term{morphism}{态射}
\[
\varepsilon_*(F)\longrightarrow\varepsilon_{Z^\bullet*}(F)^\circ
\longrightarrow\varepsilon_{Z^\bullet*}(F)^1 \tag{2.3.4.2}
\]
为零。

于是存在\Term{triangulated-functor}{三角函子}的一个\Term{canonical}{典范}\Term{morphism}{态射}
\[
Q\circ K^+(\varepsilon_*)\xrightarrow{j}
Q\circ(\ )_s\circ K^+(\varepsilon_{Z^\bullet*})
=R^+\varepsilon_{Z^\bullet*}\circ Q. \tag{2.3.4.3}
\]
\begin{SGAresult}{命题}{2.3.5}
保持 (2.3.3)、(2.3.4) 的条件和记号，并进一步假定\Term{augmentation}{增广}\Term{complex}{复形}
\[
Z^n\longrightarrow Z^{n-1}\cdots Z^1\longrightarrow Z^0\xrightarrow{t}\mathbf Z
\]
给出 $\mathbf Z$ 的一个\Term{resolution}{解消}。

则对任一由 $I_{(S\times D,\mathcal O_S)}$ 的\Term{object}{对象}组成的\Term{complex}{复形} $F^\bullet$（参见 1.3.10），
\[
j(F^\bullet):G\circ K^+(\varepsilon_*)(F^\bullet)
\longrightarrow(R^+\varepsilon_{Z^\bullet*}\circ Q)(F^\bullet)
\]
是\Term{isomorphism}{同构}。
\end{SGAresult}
\begin{proof}
设 $i$ 是 $D$ 的一个\Term{object}{对象}，$Q_i$ 是 $(S,\mathcal O_S)$ 上的\Term{module}{模}；来证明 $j(e_{i*}(Q_i))$ 是\Term{isomorphism}{同构}。也就是说，要证明列
\[
0\longrightarrow Q_i\longrightarrow\varepsilon_{Z^\bullet*}(e_{i*}(Q_i))^\circ
\longrightarrow\varepsilon_{Z^\bullet*}(e_{i*}(Q_i))^1\longrightarrow\cdots \tag{2.3.5.1}
\]
\Term{exact}{正合}。

为此，只需证明对 $\operatorname{Mod}(S,\mathcal O_S)$ 的每个\Term{object}{对象} $X$，下列序列\Term{exact}{正合}：
\[
\begin{aligned}
0&\longrightarrow\operatorname{Hom}(X,Q_i)
\longrightarrow\operatorname{Hom}(X,\varepsilon_{Z^\bullet*}(e_{i*}(Q_i))^\circ)\\
&\longrightarrow\operatorname{Hom}(X,\varepsilon_{Z^\bullet*}(e_{i*}(Q_i))^1)
\longrightarrow\cdots. 
\end{aligned}\tag{2.3.5.2}
\]
而直接计算表明，(2.3.5.2) 化为
\[
\begin{aligned}
0&\longrightarrow\operatorname{Hom}(\mathbf Z,\operatorname{Hom}(X,Q_i))\\
&\longrightarrow\operatorname{Hom}(\operatorname{Hom}_{\operatorname{Add}(D)}(Z^\circ,i),\operatorname{Hom}(X,Q_i))\\
&\longrightarrow\operatorname{Hom}(\operatorname{Hom}_{\operatorname{Add}(D)}(Z^1,i),\operatorname{Hom}(X,Q_i))\longrightarrow\cdots.
\end{aligned}\tag{2.3.5.3}
\]
若对 $D$ 的每个\Term{object}{对象} $i$，$Q_i$ 是一个\Term{totally-acyclic}{全无上同调}\Term{module}{模}，则对\Term{complex}{复形} (2.3.5.1) 应用\Term{functor}{函子} $\prod_{i\in\operatorname{Ob}(D)}$，由 (1.3.10.1) 可知，仍得到一个\Term{acyclic}{无上同调}\Term{complex}{复形}。因此 $j(\prod_{i\in\operatorname{Ob}(D)}(e_{i*}(Q_i)))$ 是\Term{isomorphism}{同构}。

\SourcePage{88}
由此可推出，当每个 $n$ 都满足 $F^n\in I_{(S\times D,\mathcal O_S)}$ 时，$j(F^\bullet)$ 是\Term{isomorphism}{同构}；这一步留给读者（例如可使用\Term{spectral-sequence}{谱序列}）。
\end{proof}
\begin{SGAresult}{推论}{2.3.6}
在 (2.3.5) 的条件下，\Term{morphism}{态射} $j$ 将 $R^+\varepsilon_{Z^\bullet*}$ 定义为 $\varepsilon_*$ 的\Term{derived-functor}{导出函子}。
\end{SGAresult}
由 (1.3.10) 和 (2.3.5)，$K^+((S\times D),\mathcal O_S)$ 中由 $I_{(S\times D,\mathcal O_S)}$ 的\Term{object}{对象}组成的\Term{complex}{复形}所定义的子\Term{category}{范畴}，对于\Term{functor}{函子} $K^+(\varepsilon_*)$ 满足 [[7], I] 中定理 (5.1) 的条件：推论立即由此得到。

\begin{SGAresult}{推论}{2.3.7}
在 (2.3.5) 的条件下，设 $(E,A)$ 是\Term{ringed}{赋环} $D$-\Term{topos}{拓扑斯}，$\theta:(E,A)\to(S\times D,\mathcal O_S)$ 是一个\Term{augmentation}{增广}；存在\Term{canonical}{典范}\Term{isomorphism}{同构}
\[
R^+\bar\theta_*\xrightarrow{\sim}R^+\Gamma(\theta_*)\circ R^+\varepsilon_{Z^\bullet*}.
\]
\end{SGAresult}
这由 [[7] I 5.4] 得到。

\begin{SGAresult}{注记}{2.3.8}
我们事先就知道 $R^+\varepsilon_*$ 存在，并且有\Term{isomorphism}{同构}
\[
R^+\bar\theta_*\xrightarrow{\sim}R^+\Gamma(\theta_*)\circ R^+\varepsilon_*,
\]
因为 $\Gamma(\theta_*)$ 和 $\varepsilon_*$ 是由\Term{ringed-topos}{赋环拓扑斯}\Term{morphism}{态射}诱导的（参见 (1.3.6) 和 (2.1.1)）。不过，前面的计算将十分有用，且可适用于其他背景，因为在实际中，(2.3.5) 的条件总会满足。
\end{SGAresult}
现在给出可以应用 (2.3.5) 的基本例子。

\textbf{2.3.9.} 在 $\operatorname{Add}(\Delta)$ 中置 $Z^n=[n]$，并由公式
\[
d^n=\sum_{i=0}^{n+1}(-1)^i\partial^i \tag{2.3.9.1}
\]
定义 $d^n:Z^n\to Z^{n+1}$，其中 $\partial^i:[n]\to[n+1]$ 是第 $i$ 个面［参见 [3] II 2］。取 $t:Z^\circ\to\mathbf Z$ 为显然的自然\Term{augmentation}{增广}。

(2.3.4) 的条件显然满足；(2.3.5) 的条件由 ([5] I (3.7.4)) 得到。

\textbf{2.3.10.} 在 $\operatorname{Add}(\Delta\times\Delta)$ 中，考虑下列\Term{double-complex}{双复形}所关联的\Term{simple-complex}{单复形}：
\[
\begin{tikzcd}
([n],[m+1])\arrow[r,"{s_1^{n,m+1}}"]&([n+1],[m+1])\\
([n],[m])\arrow[u,"{s_2^{n,m}}"]\arrow[r,"{s_1^{n,m}}"']&([n+1],[m])\arrow[u,"{s_2^{n+1,m}}"']
\end{tikzcd}\tag{2.3.10.1}
\]
\SourcePage{89}
其中
\[
\begin{cases}
S_1^{n,m}=\displaystyle\sum_{i=0}^{n+1}(-1)^i(\partial^i,\operatorname{id}_{[m]}),\\[5pt]
S_2^{n,m}=\displaystyle\sum_{i=0}^{m+1}(-1)^j(\operatorname{id}_{[n]},\partial^j).
\end{cases}\tag{2.3.10.2}
\]
利用 (2.3.9)，可验证如此定义的\Term{complex}{复形}连同显然的自然\Term{augmentation}{增广}满足 (2.3.5) 的条件；验证留给读者。

\textbf{2.3.11.} \Term{multisimplicial}{多重单纯}情形（用 $\underbrace{\Delta\times\Delta\times\cdots\times\Delta}_{p\text{ 次}}$）可作类似处理。

\textbf{2.4. \Term{relative-cohomology}{相对上同调}\Term{descent}{下降}。}
\textbf{2.4.0.} 设 $\mathscr E\to B$ 是以\Term{topos}{拓扑斯}的\Term{duality}{对偶}为\Term{fiber}{纤维}的\Term{bifibred-category}{双纤维范畴}，$\mathcal O$ 是 $(\mathscr E^\circ)$ 的一个环。

这些数据\Term{canonical}{典范}地定义一个 $B^\circ$ 上以\Term{abelian-category}{阿贝尔范畴}为\Term{fiber}{纤维}的\Term{kan-fibration}{纤维化}且余\Term{kan-fibration}{纤维化}\Term{category}{范畴}，记作 $\operatorname{Mod}(\mathscr E^\circ,\mathcal O)$（参见 (1.3.3) 和 (1.3.4)）；下文中的“\Term{direct-image}{直接像}”和“\Term{inverse-image}{逆像}”\Term{functor}{函子}总是相对于\Term{kan-fibration}{纤维化}且\Term{cofibred-functor}{余纤维化函子} $\operatorname{Mod}(\mathscr E^\circ,\mathcal O)\to B^\circ$ 而言。

\begin{SGAresult}{定义}{2.4.0.1}
若 $B$ 中的\Term{morphism}{态射} $f:T\to S$ 满足
\[
f_*:\operatorname{Mod}(\mathscr E_S^\circ,\mathcal O_S)
\longrightarrow\operatorname{Mod}(\mathscr E_T^\circ,\mathcal O_T)
\]
是\Term{exact}{正合}\Term{functor}{函子}，就称 $f$（相对于 $(\mathscr E,\mathcal O)$）\Term{flat}{平坦}。
\end{SGAresult}
\textbf{2.4.1.} 设 $G$ 是 $\operatorname{Mod}(\mathscr E^\circ,\mathcal O)$ 的一个\Term{kan-fibration}{纤维化}且余\Term{kan-fibration}{纤维化}子\Term{category}{范畴}，并且对 $B$ 的每个\Term{object}{对象} $S$，$G_S$ 是 $(\operatorname{Mod}(\mathscr E^\circ,\mathcal O))_S$ 的一个\Term{thick-subcategory}{厚子范畴}（参见 Tohokû (1.11)）。读者可在第4节找到这种情形的例子。

若 $S$ 是 $B$ 的\Term{object}{对象}，以 $D^+(S)$ 表示 $\operatorname{Mod}(\mathscr E_S^\circ,\mathcal O_S)$ 的\Term{derived-category}{导出范畴}；按照 [[7] (I §4)] 引入\Term{category}{范畴} $D^+_{G_S}(S)$，它等同于 $D^+(S)$ 的\Term{full-subcategory}{满子范畴}，由满足对每个 $n$ 都有 $H^n(X^\bullet)\in G_S$ 的\Term{object}{对象} $X^\bullet$ 组成。

始终假定满足下列条件。

\textbf{2.4.1.1.} 对 $B$ 中的任一\Term{morphism}{态射} $h:T\to S$，\Term{functor}{函子} $R^+(h^*):D^+(T)\to D^+(S)$ 满足：若 $F\in G_T$，则 $R^+(h^*)(F)\in D^+_{G_S}(S)$。

由 [[7] (I (7.3) (ii))]，此条件蕴含
\[
R^+(h^*)(D^+_{G_T}(T))\subset D^+_{G_S}(S).
\]
\textbf{2.4.2.} 设 $D$ 是 $\mathcal U$-\Term{small-category}{小范畴}：下文始终假定给定了 $\operatorname{Add}(D)$ 中一个满足 (2.3.5) 条件的\Term{complex}{复形} $Z^\bullet$（在应用中，实际上会有 $D=\Delta$ 或 $D=\Delta\times\Delta$）。

若 $X:D^\circ\to B$ 是 $B$ 的一个 $D$-\Term{object}{对象}，以 $\operatorname{Mod}_G((\underline X),\mathcal O)$ 表示 $\operatorname{Mod}((\underline X),\mathcal O)$ 的\Term{full-subcategory}{满子范畴}，由满足对 $D$ 的每个\Term{object}{对象} $i$ 都有 $e_i^*(F)\in G_i\;(=G_{X_i})$ 的\Term{object}{对象} $F$ 组成；$\operatorname{Mod}_G((\underline X),\mathcal O)$ 是 $\operatorname{Mod}((\underline X),\mathcal O)$ 的\Term{thick-subcategory}{厚子范畴}。以 $K^+_G((\underline X),\mathcal O)$（相应地，$D^+_G((\underline X),\mathcal O)$）表示 $K^+((\underline X),\mathcal O)$
\SourcePage{90}
（相应地，$D^+((\underline X),\mathcal O)$）的\Term{full-subcategory}{满子范畴}，由对每个 $n$ 都满足 $H^n(X^\bullet)\in\operatorname{Mod}_G((\underline X),\mathcal O)$ 的\Term{object}{对象} $X^\bullet$ 组成。

\textbf{2.4.3.} 设 $X,X'$ 是 $B$ 的两个 $D$-\Term{object}{对象}，$\alpha:X\to X'$ 是一个\Term{functor}{函子}\Term{morphism}{态射}。回顾（参见 (1.3.4)），我们仍以 $\alpha:(\underline X,\mathcal O)\to(\underline{X'},\mathcal O)$ 表示相应的\Term{ringed}{赋环} $D$-\Term{topos}{拓扑斯}\Term{morphism}{态射}。

\begin{SGAresult}{定义}{2.4.3.1}
若相应的\Term{ringed}{赋环} $D$-\Term{topos}{拓扑斯}\Term{morphism}{态射}在 (1.3.6.1) 的意义下\Term{flat}{平坦}，就称 $\alpha:X\to X'$ \Term{flat}{平坦}。
\end{SGAresult}
于是有下面的显然引理：
\begin{SGAresult}{引理}{2.4.3.2}
$\alpha:X\to X'$ \Term{flat}{平坦}，当且仅当对 $D$ 的每个\Term{object}{对象} $i$，$\alpha_i:X_i\to X'_i$ 都是\Term{flat}{平坦}\Term{morphism}{态射}（参见 (2.4.0.1)）。
\end{SGAresult}
\textbf{2.4.4.} 一个\Term{morphism}{态射} $\alpha:X\to X'$ 定义一个\Term{functor}{函子}
\[
K^+(\Gamma(\alpha_*)):K^+((\underline X),\mathcal O)
\longrightarrow K^+((\underline{X'}),\mathcal O). \tag{2.4.4.1}
\]
\begin{SGAresult}{引理}{2.4.4.2}
\Term{functor}{函子}
\[
K^+(\Gamma(\alpha_*))\big|_{K^+_G((\underline X),\mathcal O)}:
K^+_G((\underline X),\mathcal O)\longrightarrow K^+((\underline{X'}),\mathcal O)
\]
有一个右\Term{derived-functor}{导出函子}，记作 $R^+_G\Gamma(\alpha_*)$，并且\Term{canonical}{典范}\Term{morphism}{态射}
\[
R^+_G\Gamma(\alpha_*)\longrightarrow
R^+\Gamma(\alpha_*)\big|_{D^+_G((\underline X),\mathcal O)}
\]
是一个\Term{morphism}{态射}。

此外，$R^+_G(D^+_G((\underline X),\mathcal O))\subset D^+_G((\underline{X'}),\mathcal O)$。
\end{SGAresult}
第一部分由 (1.3.10) 和 [[7] (I (5.2))] 得到。为验证最后一个断言，由 [[7] (I (7.3) (ii))]，只须证明，当 $F\in\operatorname{Mod}_G((\underline X),\mathcal O)$ 时，$R^+\Gamma(\alpha_*)(F)\in D^+_G((\underline{X'}),\mathcal O)$。这由 (1.3.1.1) 给出的 $R^+\Gamma(\alpha_*)$ 的显式计算和假设 (2.4.1.1) 得到。

\textbf{2.4.5.} 设 $S$ 是 $B$ 的\Term{object}{对象}：\Term{ringed}{赋环} $D$-\Term{topos}{拓扑斯} $(\underline{C_S^D},\mathcal O)$（参见 (1.2.6)）等同于常值\Term{ringed}{赋环} $D$-\Term{topos}{拓扑斯} $(\mathscr E_S^\circ\times D,\mathcal O_S)$；由 [[7] (I (7.3) (ii))] 得到下列引理：
\begin{SGAresult}{引理}{2.4.5.1}
\Term{functor}{函子} $R^+\varepsilon_{Z^\bullet*}|_{D^+_G((\underline{C_S^D}),\mathcal O)}$，记作 $R^+_G\varepsilon_{Z^\bullet*}$，取值于 $D^+_{G_S}(S)$，并等同于\Term{functor}{函子}
\[
K^+(\varepsilon_*)\big|_{K^+_G((\underline{C_S^D}),\mathcal O)}:
K^+_G((\underline{C_S^D}),\mathcal O)\longrightarrow K^+(S)
\]
的一个右\Term{derived-functor}{导出函子}。
\end{SGAresult}
\begin{SGAresult}{命题}{2.4.6}
设 $\theta:X\to C_S^D$ 是 $B$ 的一个\Term{augmentation}{增广} $D$-\Term{object}{对象}：\Term{functor}{函子} $K^+(\bar\theta_*)|_{K^+_G((\underline X),\mathcal O)}$ 有一个右\Term{derived-functor}{导出函子}，记作 $R^+\bar{\mathcal O}_*$，其值落在 $D^+_{G_S}(S)$ 中。此外，存在\Term{canonical}{典范}\Term{isomorphism}{同构}
\[
R^+_G\bar{\mathcal O}_*\xrightarrow{\sim}
R^+_G\varepsilon_{Z^\bullet*}\circ R^+_G\Gamma(\theta_*).
\]
\end{SGAresult}
\SourcePage{91}
根据前述结果作出的证明留给读者。

现在可以给出有关\Term{cohomological-descent}{上同调下降}的定义。

\textbf{2.4.7.} 设 $\theta:X\to C_S^D$ 是 $B$ 的一个\Term{augmentation}{增广} $D$-\Term{object}{对象}。假定 $\theta$ \Term{flat}{平坦}，因而\Term{functor}{函子} $L^+\bar\theta^*$ 在 $D^+_{G_S}(S)$ 上的\Term{restriction}{限制}取值于 $D^+_G((\underline X),\mathcal O)$；以 $L^+_G\bar\theta^*$ 表示这一\Term{restriction}{限制}。采用这些记号：
\begin{SGAresult}{引理}{2.4.7.1}
存在两个\Term{functor}{函子}\Term{morphism}{态射}
\[
\begin{aligned}
\alpha &: \operatorname{id}_{D^+_{G_S}(S)}\longrightarrow R^+_G\bar\theta_*\circ L^+_G\bar\theta^*,\\
\beta &: L^+_G\bar\theta^*\circ R^+_G\bar\theta_*\longrightarrow\operatorname{id}_{D^+_G((\underline X),\mathcal O)},
\end{aligned}
\]
使这两个\Term{functor}{函子}构成\Term{adjunction}{伴随}。
\end{SGAresult}
这立即由 (2.2.3.1) 和 (2.4.6) 得到。

\begin{SGAresult}{定义}{2.4.8}
若满足下列条件，就称 $\theta$ 为相对于 $G$ 的 $1$-\Term{cohomological-descent}{上同调下降}\Term{augmentation}{增广}：

$1^\circ$）$\theta$ \Term{flat}{平坦}；

$2^\circ$）\Term{functor}{函子} $L^+_G\bar\theta^*$ \Term{fully-faithful}{充分忠实}。
\end{SGAresult}
\begin{SGAresult}{注记}{2.4.9}
前一定义中的条件 $2^\circ$ 也可以表述为：(2.4.7.1) 中的\Term{morphism}{态射} $\alpha$ 是\Term{isomorphism}{同构}。
\end{SGAresult}
\begin{SGAresult}{定义}{2.4.10}
设 $\theta:X\to C_S^D$ 是 $B$ 的一个\Term{augmentation}{增广} $D$-\Term{object}{对象}。若\Term{functor}{函子}
\[
\Gamma(\theta^*)\circ\varepsilon^*:G_S\longrightarrow\operatorname{Cocart}(D,G)
\]
是\Term{category}{范畴}\Term{equivalence}{等价}，就称 $\theta$ 为相对于 $G$ 的\Term{effective-descent}{有效下降}\Term{augmentation}{增广}。
\end{SGAresult}
\begin{SGAresult}{定义}{2.4.11}
若 $\theta$ 同时是相对于 $G$ 的 $1$-\Term{cohomological-descent}{上同调下降}\Term{augmentation}{增广}和\Term{effective-descent}{有效下降}\Term{augmentation}{增广}，就称 $\theta$ 为相对于 $G$ 的 $2$-\Term{cohomological-descent}{上同调下降}\Term{augmentation}{增广}（或\Term{effective-cohomological-descent}{有效上同调下降}\Term{augmentation}{增广}）。
\end{SGAresult}
由 (2.2.7) 的证明得到下述结果：
\begin{SGAresult}{定理}{2.4.12}
设 $\theta:X\to C_S^D$ 是相对于 $G$ 的 $2$-\Term{cohomological-descent}{上同调下降}\Term{augmentation}{增广}。则 $L^+_G\bar\theta^*$ 的\Term{essential-image}{本质像}是 $D^+_G((\underline X),\mathcal O)$ 的\Term{full-subcategory}{满子范畴}，由满足下列条件的\Term{complex}{复形} $F^\bullet$ 组成：对每个 $n$，$H^n(F^\bullet)$ 是 $G$ 的\Term{cocartesian}{余笛卡尔}\Term{section}{截面}。
\end{SGAresult}
下列命题使我们能够用 ([4] 6) 的语言重述定义 (2.4.10)，其验证留给读者：
\begin{SGAresult}{命题}{2.4.13}
$\theta:X\to C_S^D$ 是相对于 $G$ 的\Term{effective-descent}{有效下降}\Term{augmentation}{增广}，当且仅当 $\theta$ 定义的 $B$ 上的\Term{functor}{函子} $D^\circ\to B/S$，在 $G^\circ$ 于 $B/S$ 上的\Term{cartesian}{笛卡尔}\Term{section}{截面}\Term{category}{范畴}和 $G^\circ$ 于 $D^\circ$ 上的\Term{cartesian}{笛卡尔}\Term{section}{截面}\Term{category}{范畴}之间诱导一个\Term{equivalence}{等价}。
\end{SGAresult}

\SourcePage{92}
\begin{SGAresult}{推论}{2.4.14}
假定 $B$ 中的\Term{finite}{有限}\Term{fiber-product}{纤维积}\Term{representable-functor}{可表示}，设 $f:R\to S$ 是一个\Term{morphism}{态射}：\Term{augmentation}{增广}\Term{semisimplicial-object}{半单纯对象} $[R|_fS]$（参见 (1.2.7)）是相对于 $G$ 的\Term{effective-descent}{有效下降}，当且仅当 $f$ 对于 $B$ 上的\Term{fibred-category}{纤维范畴} $G^\circ$ 是一个\Term{effective-descent}{有效下降}\Term{morphism}{态射}。
\end{SGAresult}
这由 (2.4.13) 和 ([4] 9) 得到。

\begin{SGAresult}{定义}{2.4.15}
假定 $B$ 中的\Term{finite}{有限}\Term{fiber-product}{纤维积}\Term{representable-functor}{可表示}，设 $f:R\to S$ 是一个\Term{morphism}{态射}。若\Term{augmentation}{增广}\Term{semisimplicial-object}{半单纯对象} $[R|_fS]$ 是相对于 $G$ 的 $i$-\Term{cohomological-descent}{上同调下降}，就称 $f$ 是相对于 $G$ 的 $i$-\Term{cohomological-descent}{上同调下降}\Term{morphism}{态射}（$i=1,2$）。
\end{SGAresult}
\begin{SGAresult}{定义}{2.4.16}
假定 $B$ 中的\Term{finite}{有限}\Term{fiber-product}{纤维积}\Term{representable-functor}{可表示}，设 $\theta:X\to C_S^D$ 是一个\Term{augmentation}{增广}：若对每个\Term{morphism}{态射} $T\to S$，由\Term{base-change}{基变换}得到的\Term{augmentation}{增广} $X_T\xrightarrow{\theta_T}C_T^D$ 都是相对于 $G$ 的 $i$-\Term{cohomological-descent}{上同调下降}\Term{augmentation}{增广}，就称 $\theta$ 为相对于 $G$ 的\Term{universal}{普适} $i$-\Term{cohomological-descent}{上同调下降}\Term{augmentation}{增广}（$i=1,2$）。
\end{SGAresult}
相对于 $G$ 的\Term{universal}{普适} $i$-\Term{cohomological-descent}{上同调下降}\Term{morphism}{态射}的概念，立即由前两个定义得到。

我们常用术语“$G$-$i$-\Term{cohomological-descent}{上同调下降}\Term{augmentation}{增广}”代替“相对于 $G$ 的 $i$-\Term{cohomological-descent}{上同调下降}\Term{augmentation}{增广}”。

\textbf{2.5. \Term{descent-spectral-sequence}{下降谱序列}。}
\textbf{2.5.1.} 设 $A$ 是\Term{abelian-category}{阿贝尔范畴}。以 $F(A)$ 表示如下\Term{category}{范畴}：其\Term{object}{对象}是配有\Term{discrete}{离散}且余\Term{discrete}{离散}\Term{filtration}{滤过}的 $A$ 的\Term{object}{对象}，其\Term{morphism}{态射}是 $A$ 中的\Term{filtration}{滤过}\Term{morphism}{态射}；\Term{category}{范畴} $F(A)$ 是\Term{additive-category}{加性范畴}（但并非\Term{abelian-category}{阿贝尔范畴}）。

$A$ 中逐次次数的\Term{filtration}{滤过}均\Term{discrete}{离散}且余\Term{discrete}{离散}的\Term{filtration}{滤过}\Term{complex}{复形}所成的\Term{category}{范畴} $K(F(A))$ 是\Term{triangulated-category}{三角范畴}，\Term{canonical}{典范}\Term{functor}{函子}
\[
\operatorname{Gr}_n:K(F(A))\longrightarrow K(A) \tag{2.5.1.1}
\]
是\Term{triangulated-functor}{三角函子}。因此 $K(F(A))$ 中使得对每个 $n$，$\operatorname{Gr}_n(f)$ 都是\Term{quasi-isomorphism}{拟同构}的\Term{morphism}{态射} $f$ 所成的\Term{set}{集合} $\Sigma$ 是一个\Term{saturated-multiplicative-system}{饱和乘法系}［参见 [9] §2，第1号］。

将这个系中的\Term{arrow}{箭头}逆化，得到一个新的\Term{triangulated-category}{三角范畴}，记作 $DF(A)$，而\Term{functor}{函子} $\operatorname{Gr}_n$ 延拓为\Term{functor}{函子}
\[
\operatorname{Gr}_n:DF(A)\longrightarrow D(A). \tag{2.5.1.2}
\]
此后只考虑\Term{bounded-below}{下有界}\Term{complex}{复形}；自然引入记号 $K^+(F(A))$ 和 $D^+F(A)$。

\SourcePage{93}
\textbf{2.5.2.} 设 $B$ 是 $A$ 的\Term{thick-subcategory}{厚子范畴}，以 $K^+_B(F(A))$ 表示 $K^+(F(A))$ 的\Term{full-subcategory}{满子范畴}，其\Term{object}{对象}是满足对每对整数 $(i,j)$ 都有 $H^i(\operatorname{Gr}_j(X^\bullet))\in B$ 的\Term{complex}{复形} $X^\bullet$；$K^+_B(F(A))$ 是 $K(F(A))$ 的一个\Term{localization}{局部化}\Term{triangulated-subcategory}{三角子范畴}［参见 [7] (I §5)］。以 $D^+_BF(A)$ 表示 $DF(A)$ 的\Term{full-subcategory}{满子范畴}，其\Term{object}{对象}是满足对每对 $(i,j)$ 都有 $H^j(\operatorname{Gr}_j(X^\bullet))\in B$ 的\Term{bounded-below}{下有界}\Term{complex}{复形} $X^\bullet$。根据同上，$D^+_BF(A)$ 等同于\Term{category-of-fractions}{分式范畴} $K^+_B(F(A))_{\Sigma\cap K^+_B(F(A))}$。

“遗忘\Term{filtration}{滤过}”\Term{functor}{函子}
\[
\iota:K^+(F(A))\longrightarrow K^+(A)
\]
是\Term{triangulated-functor}{三角函子}。此外，由 [[5] I (4.7)]，若 $f\in\Sigma$，则 $\iota(f)$ 是\Term{quasi-isomorphism}{拟同构}，且有 $\iota(K^+_B(F(A)))\subset K^+_B(A)$。于是\Term{functor}{函子} $\iota$ 延拓为\Term{triangulated-functor}{三角函子}
\[
\iota:D^+F(A)\longrightarrow D^+(A), \tag{2.5.2.1}
\]
满足 $\iota(D^+_BF(A))\subset D^+_B(A)$。

最后注意，存在一个\Term{canonical}{典范}\Term{spectral-functor}{谱函子} $D^+_BF(A)\to B$，它收敛到 $H^n\circ\iota$，而其 $E_1$ 项由 $H^{p+q}\circ\operatorname{Gr}_p$ 给出［参见 [5] I (4.2)］。

现在回到 (2.3) 的记号，并假定 $D=\Delta$；以 $Z^\bullet$ 表示由 (2.3.9.1) 定义的 $\operatorname{Add}(\Delta)$ 中的\Term{complex}{复形}。

\begin{SGAresult}{命题}{2.5.3}
设 $(S,\mathcal O_S)$ 是\Term{ringed-topos}{赋环拓扑斯}。\Term{functor}{函子}
\[
R^+\varepsilon_{Z^\bullet*}:D^+((S\times D),\mathcal O_S)\longrightarrow D^+(S,\mathcal O_S)
\]
有一个\Term{canonical}{典范}\Term{resolution}{分解}
\[
D^+(\Gamma(S\times D),\mathcal O_S)
\xrightarrow{\mathrm F R^+\varepsilon_{Z^\bullet*}}
D^+F(S,\mathcal O_S)\xrightarrow{\iota}D^+(S,\mathcal O_S).
\]
此外，对每个整数 $i$，下图交换：
\[
\begin{tikzcd}[column sep=large]
D^+((S\times D),\mathcal O_S)\arrow[r,"R^+(e_i^*)"]\arrow[d,"\mathrm F R^+\varepsilon_{Z^\bullet*}"']&D^+(S,\mathcal O_S)\\
D^+F(S,\mathcal O_S)\arrow[ur,"\operatorname{Gr}_i"']&
\end{tikzcd}
\]
\end{SGAresult}
\begin{proof}
考虑\Term{functor}{函子} $C^+(C^+(S,\mathcal O_S))\xrightarrow{(\ )_s}C^+(S,\mathcal O_S)$；若 $X^{\bullet\bullet}$ 是 $C^+(C^+(S,\mathcal O_S))$ 的\Term{object}{对象}，则为 $(X^{\bullet\bullet})_s$ 配以第二\Term{canonical}{典范}\Term{filtration}{滤过}［参见 [5] I 4.8］，从而得到\Term{resolution}{分解}
\[
(\ )_s:C^+(C^+(S,\mathcal O_S))\longrightarrow
C^+(F(S,\mathcal O_S))\xrightarrow{\iota}C^+(S,\mathcal O_S). \tag{2.5.3.1}
\]
\SourcePage{94}
它\Term{descent}{下降}到\Term{category}{范畴} $K^+$：
\[
(\ )_s:K^+(C^+(S,\mathcal O_S))\longrightarrow
K^+(F(S,\mathcal O_S))\xrightarrow{\iota}K^+(S,\mathcal O_S), \tag{2.5.3.2}
\]
因为 $C^+(C^+(S,\mathcal O_S))$ 中的\Term{homotopy}{同伦}诱导一个\Term{filtration}{滤过}\Term{homotopy}{同伦}。

于是有下图：
\[
\begin{tikzcd}[column sep=small,row sep=large]
D^+((S\times\Delta),\mathcal O_S)\arrow[rr,"R\varepsilon_{Z^\bullet*}"]\arrow[dr,dashed]&&D^+(S,\mathcal O_S)\\
&D^+F(S,\mathcal O_S)\arrow[ur,"\iota"']&\\
&K^+(F(S,\mathcal O_S))\arrow[u]\arrow[dr,"\iota"]&\\
K^+((S\times\Delta),\mathcal O_S)\arrow[uuu]\arrow[ur]\arrow[rr,"{(\ )_s,\,K^+(\varepsilon_{Z^\bullet*})}"']&&K^+(S,\mathcal O_S)\arrow[uuu]
\end{tikzcd}\tag{2.5.3.3}
\]
可验证，恰有一个\Term{functor}{函子}
\[
\mathrm F R^+\varepsilon_{Z^\bullet*}:D^+((S\times\Delta),\mathcal O_S)\longrightarrow D^+F(S,\mathcal O_S)
\]
使 (2.5.3.3) 交换。

命题的其余部分是显然的。
\end{proof}

\begin{SGAresult}{命题}{2.5.4}
设 $(E,A)$ 是一个\Term{ringed}{赋环}半\Term{simplicial-topos}{单纯拓扑斯}，$\theta:(E,A)\to(S\times\Delta,\mathcal O_S)$ 是一个\Term{augmentation}{增广}。对每个整数 $i$，以 $\theta_i:(E_i,A_i)\to(S,\mathcal O_S)$ 表示 $\theta$ 在 $i$ 上方诱导的\Term{ringed-topos}{赋环拓扑斯}\Term{morphism}{态射}。

\Term{functor}{函子} $R^+\bar\theta_*:D^+((E),A)\to D^+(S,\mathcal O_S)$ 有一个\Term{canonical}{典范}\Term{resolution}{分解}
\[
D^+((E),A)\xrightarrow{\mathrm F R^+\bar\theta_*}D^+F(S,\mathcal O_S)
\xrightarrow{\iota}D^+(S,\mathcal O_S),
\]
使得对每个整数 $i$，\Term{canonical}{典范}图
\[
\begin{tikzcd}[column sep=large]
D^+((E),A)\arrow[r,"R^+e_i^*"]\arrow[d,"\mathrm F R^+\bar\theta_*"']&D^+(E_i,A_i)\arrow[d,"R^+\theta_{i*}"]\\
D^+F(S,\mathcal O_S)\arrow[r,"\operatorname{Gr}_i"']&D^+(S,\mathcal O_S)
\end{tikzcd}
\]
本质交换。
\end{SGAresult}
这由前面的结果和 (1.3.1.2) 得到。

\SourcePage{95}
\begin{SGAresult}{命题}{2.5.5}
设 $(E,A)$ 是\Term{ringed}{赋环}半\Term{simplicial-topos}{单纯拓扑斯}，$\theta:(E,A)\to(S\times\Delta,\mathcal O_S)$ 是一个 $1$-\Term{cohomological-descent}{上同调下降}\Term{augmentation}{增广}。设 $H:(S,\mathcal O_S)\to(R,\mathcal O_R)$ 是\Term{ringed-topos}{赋环拓扑斯}\Term{morphism}{态射}。则存在从 $D^+(S,\mathcal O_S)$ 到 $\operatorname{Mod}(R,\mathcal O_R)$ 的一个\Term{spectral-functor}{谱函子}
\[
E_1^{pq}=R^q(H\circ\theta_p)_*\circ L^+\theta_p^*
\longrightarrow R^{p+q}H_*,
\]
称为\Term{descent-spectral-functor}{下降谱函子}。
\end{SGAresult}
\begin{proof}
设 $H\circ\theta:(E,A)\to(R\times\Delta,\mathcal O_R)$ 是由 $\theta$ 与 $H$ 复合而\Term{canonical}{典范}得到的\Term{augmentation}{增广}；有
\[
\overline{(H\circ\theta)}_*=H_*\circ\bar\theta_*, \tag{2.5.5.1}
\]
所以图
\[
\begin{tikzcd}[column sep=large]
D^+(S,\mathcal O_S)\arrow[r,"L^+(\bar\theta^*)"]\arrow[dr,"R^+H_*"']&D^+((E),A)\arrow[d,"R^+\overline{(H\circ\theta)}_*"]\\
&D^+(R,\mathcal O_R)
\end{tikzcd}\tag{2.5.5.2}
\]
本质交换。
\[
\left(R^+\overline{(H\circ\theta)}_*\xrightarrow{\sim}R^+H_*\circ R^+\bar\theta_*
\quad\text{且}\quad
\operatorname{id}_{D^+(S,\mathcal O_S)}\xrightarrow{\sim}R^+\bar\theta_*\circ L^+\bar\theta_*\right).
\]
由 (2.5.4)，对每个 $i$ 有本质\Term{commutative-diagram}{交换图}：
\[
\begin{tikzcd}[column sep=small,row sep=large]
D^+(S,\mathcal O_S)\arrow[r,"L^+\bar\theta^*"]\arrow[ddr,"R^+H_*"']&D^+((E),A)\arrow[r,"R^+(e_i^*)"]\arrow[d,"\mathrm F R^+\overline{(H\circ\theta)}_*"']&D^+(E_i,A_i)\arrow[d,"{R^+(R,\theta_i)_*}"]\\
&D^+F(R,\mathcal O_R)\arrow[r,"\operatorname{Gr}_i"]\arrow[d,"\iota"]&D^+F(R,\mathcal O_R)\\
&D^+(R,\mathcal O_R)&
\end{tikzcd}\tag{2.5.5.3}
\]
所宣称的\Term{spectral-functor}{谱函子}，由 $D^+F(R,\mathcal O_R)$ 上的\Term{canonical}{典范}\Term{spectral-functor}{谱函子}与
\[
\mathrm F R^+\overline{(H\circ\theta)}_*\circ L^+(\bar\theta^*)
\]
复合得到。
\end{proof}
下列命题的验证留给读者：

\SourcePage{96}
\begin{SGAresult}{命题}{2.5.6}
采用前述记号，前一个\Term{spectral-functor}{谱函子}的 $E_2^{pq}$ 项写作
\[
\check H^p\circ R^q\Gamma((H\circ\theta)_*)\circ L^+\bar\theta^*,
\]
其中 $\check H^p:\operatorname{Mod}((R\times\Delta),\mathcal O_R)\to\operatorname{Mod}(R,\mathcal O_R)$ 是如下\Term{functor}{函子}：给每个\Term{functor}{函子} $\Delta\xrightarrow{F}\operatorname{Mod}(R,\mathcal O_R)$ 配以其关联\Term{complex}{复形}的 $p$ 次\Term{homology}{同调}\Term{object}{对象}（该\Term{complex}{复形}在 (2.3.3) 中记作 $\varepsilon_{Z^\bullet*}(F)^\bullet$）。
\end{SGAresult}
现在回到 (2.4) 引入的术语：借助 (2.5.2) 的一般性事实，在所有有意义的位置加上下标字母 $G$，便可将前面的一切讨论逐字改写。细节留给读者；特别得到：
\begin{SGAresult}{命题}{2.5.7}
设 $\mathcal O:X\to C_S^D$ 是相对于 $G$ 的 $1$-\Term{cohomological-descent}{上同调下降}\Term{augmentation}{增广}，$h:S\to R$ 是 $B$ 中的\Term{morphism}{态射}；存在从 $D^+_{G_S}(S)$ 到 $G_R$ 的\Term{spectral-functor}{谱函子}
\[
E_1^{pq}=R_G^q(h\circ\theta)_*\circ L_G^+\theta_{p*}
\longrightarrow R_G^{p+q}h_*,
\]
且有 $E_2^{pq}=\check H^p\circ R_G^q\Gamma((h\circ\theta)_*)\circ L_G^+\bar\theta^*$。
\end{SGAresult}
% 本块在原 PDF 第96页“3. Critères de descente”标题前结束。
